\chapter{The involution \texorpdfstring{$\theta$}{theta} and the fixed-point algebra \texorpdfstring{$V_\Lambda^+$}{V+}}\label{ch:04}

\providecommand{\chfourone}{\mathbf{1}}

This chapter documents the release of the LeanVOA library named in the front matter; every Lean quotation below is checked by \texttt{scripts/textbook/sync.py}
to occur verbatim (up to whitespace and comments) in the source files of that release. Checks
that were run at an earlier state of the library and not re-run for this edition are marked as such.
The Moonshine module of Frenkel, Lepowsky and Meurman is assembled from two halves,
\begin{equation*}
  \Vnat \;=\; V_\Leech^{+}\;\oplus\;(V_\Leech^{T})^{+},
\end{equation*}
an \emph{untwisted} half and a \emph{twisted} half. This chapter is entirely about the first
summand. Starting from the Leech lattice vertex operator algebra $V_\Leech$ (constructed in
Chapter~\ref{ch:03}) we lift the isometry $-1$ of the lattice to an automorphism $\theta$ of
$V_\Leech$, form the subalgebra $V_\Leech^+$ of $\theta$-fixed vectors, and establish everything the
later chapters need to know about it:
\begin{enumerate}[label=(\roman*)]
\item $\theta$ is a genuine VOA automorphism of order exactly two
  (\S\ref{ch04:sec:theta}--\S\ref{ch04:sec:aut});
\item $V_\Leech^+$ is a vertex operator algebra of central charge $24$ (\S\ref{ch04:sec:plus});
\item $(V_\Leech^+)_1 = 0$, $\dim (V_\Leech^+)_2 = 98580$, $\dim (V_\Leech^+)_3 = 10745856$, and
  more generally an exact formula for every $\dim(V_\Leech^+)_k$
  (\S\ref{ch04:sec:wtone}--\S\ref{ch04:sec:gradeddim});
\item $V_\Leech^+$ is generated, as a vertex algebra, by the vectors $e^\beta+e^{-\beta}$
  (\S\ref{ch04:sec:generation});
\item the canonical invariant bilinear form of $V_\Leech^+$ is nondegenerate on every weight space,
  proved by writing down an explicit form on all of $V_\Leech$ (\S\ref{ch04:sec:form}).
\end{enumerate}
Item (v) is the input that makes the twisted-by-twisted operator $\overline W$ of
Chapter~\ref{ch:09} nondegenerate, which in turn is what distinguishes $\Vnat$ from a degenerate
``square-zero'' impostor with the same graded dimension; we shall only point at that use here.

\paragraph{The physical picture.} A rank-$24$ lattice VOA is the chiral half of the conformal
field theory of $24$ free bosons $X^a(z)$ compactified on the torus $\mathbb R^{24}/\Leech$. Its
states are oscillator excitations $h_{a}(-n)$ of the bosons, dressed by a momentum (``charge'')
$\alpha\in\Leech$. The reflection $X\mapsto -X$ of the torus is a symmetry: it flips the sign of every
oscillator and of every momentum. The quantum operator that implements this reflection is $\theta$.
Gauging the reflection (``orbifolding'') has two effects on the space of states: one keeps only the
reflection-invariant states of the original theory --- this is $V_\Leech^+$ --- and one adds a
\emph{twisted sector} of strings closed only up to the reflection --- that is $(V_\Leech^T)^+$,
the subject of Chapters~\ref{ch:06}--\ref{ch:08}. Seen this way, this chapter computes the
\emph{projected untwisted sector}: which states survive the projection, how many there are at
each energy, how they are generated, and which inner product they carry.

\paragraph{How to read this chapter.} Every definition and theorem is presented with the
Fidelity Protocol of the book: the mathematical statement, then the Lean declaration copied
verbatim from the repository, then a fidelity check (in a box titled ``From the mathematics to the
Lean statement'') that decodes every symbol and hypothesis, then a proof on paper, then the
machine status. The files of this chapter are, in \leanpath{VOA/Moonshine/},
\leanpath{LatticeInvolution.lean}, \leanpath{PlusWeightOne.lean}, \leanpath{PlusWeightTwo.lean},
\leanpath{LatticeGradedDim.lean}, \leanpath{GenerationFock.lean}, \leanpath{Generation.lean},
\leanpath{NondegFock.lean}, \leanpath{NondegLattice.lean} and \leanpath{NondegLeech.lean}.

\begin{status}
All statements quoted in this chapter are in files imported by the umbrella
\leanpath{VOA/Moonshine.lean} (checked at an earlier commit, when this chapter was first
written, and again in the working tree of a pre-release build commit); the
aggregate \texttt{lake build} recorded by the project for that earlier state was green
(10343 jobs; not re-run for this chapter). While writing this chapter the command
\texttt{lake env lean -DmaxSynthPendingDepth=3} was run on a scratch file importing
\leanpath{VOA.Moonshine.NondegLeech} and \leanpath{VOA.Moonshine.LatticeGradedDim}, printing the
axioms of nineteen headline declarations of the chapter (\lean{thetaLat_mode},
\lean{thetaAut_sq}, \lean{thetaAut_ne_one}, \lean{leechTheta_ne_one},
\lean{mem_leechVOAPlus_iff}, \lean{leechVOAPlus_centralCharge}, \lean{leechVOAPlus_grading_one},
the four dimension theorems at weights $2$ and $3$, \lean{hasProd_fockGenFun},
\lean{hasProd_fockSignedGenFun}, \lean{two_mul_finrank_plusK_eq_signed},
\lean{leechPlus_generated}, \lean{latForm_nondeg}, \lean{latForm_Yeαvec},
\lean{leechBV_eq_plusForm}, \lean{leechBV_gradedNondeg}). Every one printed exactly
\texttt{[propext, Classical.choice, Quot.sound]}: Lean's three standard axioms, and in particular
\emph{not} \lean{VOA.pbw} (the PBW theorem of Chapter~\ref{ch:L02}; a project axiom in earlier versions of the library, now a theorem). No theorem of the chapter has an unproved datum
among its hypotheses; the only hypotheses anywhere are evenness and positive-definiteness of a
lattice, which for the Leech lattice are themselves proved theorems.
\end{status}

%======================================================================
\section{What we take from the lattice VOA}\label{ch04:sec:recall}

The lattice VOA $V_L$ is constructed in Chapter~\ref{ch:03}. To keep this chapter self-contained
we recall every definition that appears in a statement below, with its Lean source. A reader who
has just read Chapter~\ref{ch:03} may skim this section.

\subsection{Lattices}

\begin{definition}[Integral lattice]\label{ch04:def:lattice}
An \emph{integral lattice} is a free $\Z$-module $L$ of finite rank together with a symmetric
$\Z$-bilinear form $\langle\cdot,\cdot\rangle:L\times L\to\Z$. It is \emph{even} if
$\langle v,v\rangle\in 2\Z$ for all $v$, and \emph{positive-definite} if $\langle v,v\rangle>0$ for
$v\neq0$. The \emph{half-norm} of $\alpha$ is $\tfrac12\langle\alpha,\alpha\rangle$, an integer
when $L$ is even; it will be the conformal weight of $e^\alpha$.
\end{definition}

In Lean an integral lattice is the structure \lean{VOA.IntegralLattice} (carrier a free, finitely
generated $\Z$-module, symmetric form \lean{L.form}), with the predicates \lean{L.Even} and
\lean{L.PosDef} (\leanpath{VOA/Foundations/Lattices.lean}); these are quoted verbatim and decoded in
Definitions~\ref{ch02:def:voalattice}, \ref{ch02:def:voaevenposdef} and~\ref{l02:def:lattice}.
The half-norm is
\begin{leancode}
-- VOA/Examples/LatticeGenGrading.lean
noncomputable def halfNorm (α : L.carrier) : ℤ := L.form α α / 2
\end{leancode}

\begin{dictionary}
\lean{L.form α β} is $\langle\alpha,\beta\rangle$. \lean{halfNorm} uses integer division
\lean{/ 2}; this equals $\tfrac12\langle\alpha,\alpha\rangle$ exactly when $\langle\alpha,\alpha\rangle$
is even, which is why every theorem that uses \lean{halfNorm} as a conformal weight carries the
hypothesis \lean{hEven : L.Even} (the library records $2\cdot$\lean{halfNorm L α}
$=\langle\alpha,\alpha\rangle$ as \lean{two_halfNorm}, under \lean{hEven}). The \emph{rank}
\lean{rank L} is \lean{Module.finrank ℤ L.carrier}, the number of elements of a $\Z$-basis.
\end{dictionary}

\begin{physics}
An even lattice is the data of the momentum (charge) lattice of a chiral boson compactified on the torus $\mathbb R^r/L$: the evenness condition says the vertex operator $:\!e^{i\alpha\cdot X}\!:$ has integer conformal weight $\tfrac12\langle\alpha,\alpha\rangle$ (so the fields are mutually local, with integer spin), and positive-definiteness says the energy is bounded below. The Leech lattice is special in having no vectors of norm $2$, i.e.\ no extra weight-one currents, so the torus theory has only the $24$ abelian $U(1)$ currents $\partial X^a$. Heuristic, not part of the proof.
\end{physics}

The Leech lattice enters through the declaration \lean{VOA.Moonshine.leech}. It is the Leech
lattice of the vendored library in the directory \leanpath{VOA/LatticesAndCodes/}
(Chapter~\ref{ch:02}), transported into the structure above:
\begin{leancode}
-- VOA/Moonshine/LeechLatticeVOA.lean
noncomputable abbrev leech : VOA.IntegralLattice.{0} := ofEbeling Leech.lattice
theorem leech_even : leech.Even := ofEbeling_even _ Leech.lattice_isEven
theorem leech_posDef : leech.PosDef := ofEbeling_posDef _ Leech.lattice_isPosDef
theorem leech_rank : rank leech = 24 := Leech.lattice_rank
theorem leech_no_roots (v : leech.carrier) : leech.form v v ≠ 2 := Leech.lattice_no_roots v
\end{leancode}
These facts --- even, positive-definite, rank $24$, no vectors of norm $2$ --- together with the
counts $\#\{\alpha:\langle\alpha,\alpha\rangle=4\}=196560$ (\lean{Leech.repCount_four}) and
$\#\{\alpha:\langle\alpha,\alpha\rangle=6\}=16773120$ (\lean{Leech.repCount_six}) are the only
properties of $\Leech$ used in this chapter. All are proved theorems of the vendored library
(Chapter~\ref{ch:02}).

\subsection{The Fock space}

\begin{definition}[Rank-$r$ Fock space]\label{ch04:def:fock}
Let $r\ge0$. The Fock space of $r$ free bosons is the polynomial ring
\begin{equation*}
  F_r=\C\big[\,x_{(m,a)} : m\in\N,\ a\in\{0,\dots,r-1\}\,\big],
\end{equation*}
where the variable $x_{(m,a)}$ stands for the creation operator $h_a(-m-1)$ applied to the vacuum
$1$. The oscillators act by
\begin{equation*}
  h_a(-n)=\text{multiplication by }x_{(n-1,a)}\ (n\ge1),\qquad
  h_a(n)=n\,\frac{\partial}{\partial x_{(n-1,a)}}\ (n\ge1),\qquad h_a(0)=0,
\end{equation*}
so that $[h_a(m),h_b(n)]=m\,\delta_{ab}\,\delta_{m+n,0}$. The \emph{weight} of the monomial
$x^d=\prod x_{(m,a)}^{d_{(m,a)}}$ is $\sum (m+1)\,d_{(m,a)}$, its \emph{degree} is
$\sum d_{(m,a)}$ (the number of oscillators), and $(F_r)_w$ is the span of the monomials of
weight $w$.
\end{definition}

\begin{leancode}
-- VOA/Examples/HeisenbergNFock.lean
abbrev FockN (r : ℕ) : Type := MvPolynomial (ℕ × Fin r) ℂ

noncomputable def oscN (r : ℕ) (k : ℂ) (a : Fin r) (n : ℤ) : FockN r →ₗ[ℂ] FockN r :=
  if n < 0 then mulXN (-n - 1).toNat a
  else if 0 < n then (n : ℂ) • k • pderivLMN (n - 1).toNat a
  else 0

def weightFnN {r : ℕ} : ℕ × Fin r → ℕ := fun u => u.1 + 1

def gradingWN (r : ℕ) (w : ℕ) : Submodule ℂ (FockN r) :=
  weightedHomogeneousSubmodule ℂ weightFnN w

-- VOA/Examples/HeisenbergNVertex.lean
noncomputable def gradingZN (r : ℕ) (n : ℤ) : Submodule ℂ (FockN r) :=
  if 0 ≤ n then gradingWN r n.toNat else ⊥
\end{leancode}

\begin{dictionary}
\lean{MvPolynomial (ℕ × Fin r) ℂ} is Mathlib's polynomial ring over $\C$ in variables indexed by
pairs $(m,a)$ with $m\in\N$ and $a\in\{0,\dots,r-1\}$ (\lean{Fin r}). The variable
\lean{X (m,a)} is $x_{(m,a)}=h_a(-m-1)\chfourone$: the index is shifted by one, $m=0$ being the
first creation mode $h_a(-1)$. \lean{oscN r k a n} is $h_a(n)$ at level $k$; every use in this
chapter has $k=1$. \lean{mulXN i a} multiplies by \lean{X (i,a)} and \lean{pderivLMN i a} is
$\partial/\partial x_{(i,a)}$; \lean{→ₗ[ℂ]} denotes a $\C$-linear map. A monomial is recorded by its
exponent function $d$ with finite support (Mathlib's \lean{(ℕ × Fin r) →₀ ℕ}); its weight is
\lean{Finsupp.weight weightFnN d} $=\sum(m+1)d_{(m,a)}$ and its degree is \lean{d.degree}
$=\sum d_{(m,a)}$. \lean{gradingWN r w} is the subspace (\lean{Submodule}) of polynomials all of whose
monomials have weight $w$ (Mathlib's ``weighted-homogeneous'' submodule); \lean{gradingZN} extends
the index to $\Z$ by the zero subspace \lean{⊥} in negative weight.
\end{dictionary}

\begin{physics}
The Fock space is the space of oscillator excitations of $r$ free chiral bosons about the vacuum: a monomial is a state with $\deg d$ quanta, and its weight is the total energy $\sum(m+1)d_{(m,a)}$ in units where $h_a(-n)$ raises $L_0$ by $n$. The generating function of monomials by weight is the partition function of $r$ free bosons, $\prod(1-q^n)^{-r}$, computed later in this chapter. The relation $[h_a(m),h_b(n)]=m\,\delta_{ab}\delta_{m+n,0}$ is the canonical commutation relation of the current $i\partial X^a$.
\end{physics}

\subsection{The lattice VOA as a vector space}

\begin{definition}[$V_L$ as a graded vector space]\label{ch04:def:VL}
For an integral lattice $L$ of rank $r$, put $\C\{L\}=\bigoplus_{\alpha\in L}\C e^\alpha$ (a vector
space with basis $e^\alpha$; no product is needed in this chapter) and
\begin{equation*}
  V_L = F_r\otimes_{\C} \C\{L\}.
\end{equation*}
The vector $p\otimes e^\alpha$ has \emph{Fock weight} $w$ if $p\in(F_r)_w$, \emph{charge} $\alpha$,
and \emph{conformal weight} $w+\tfrac12\langle\alpha,\alpha\rangle$. The weight-$k$ space is
\begin{equation*}
  (V_L)_k=\bigoplus_{\alpha\in L}(F_r)_{k-\frac12\langle\alpha,\alpha\rangle}\otimes e^\alpha .
\end{equation*}
We write $e^\alpha$ also for the ``charged vacuum'' $1\otimes e^\alpha$; $e^0$ is the vacuum.
\end{definition}

\begin{leancode}
-- VOA/Examples/LatticeGenModule.lean
noncomputable abbrev LatModuleGen : Type :=
  HeisenbergN.FockN (rank L) ⊗[ℂ] (L.carrier →₀ ℂ)
noncomputable def eGen (α : L.carrier) : L.carrier →₀ ℂ := Finsupp.single α 1
noncomputable def latStateGen (p : HeisenbergN.FockN (rank L)) (α : L.carrier) : LatModuleGen L :=
  p ⊗ₜ[ℂ] eGen L α
noncomputable def latChargeGen (α : L.carrier) : LatModuleGen L :=
  latStateGen L (HeisenbergN.vacuumN (rank L)) α

-- VOA/Examples/LatticeGenGrading.lean
noncomputable def mkEα (α : L.carrier) : HeisenbergN.FockN (rank L) →ₗ[ℂ] LatModuleGen L :=
  (TensorProduct.mk ℂ (HeisenbergN.FockN (rank L)) (L.carrier →₀ ℂ)).flip (eGen L α)
noncomputable def bigradα (wα : ℤ × L.carrier) : Submodule ℂ (LatModuleGen L) :=
  (HeisenbergN.gradingZN (rank L) wα.1).map (mkEα L wα.2)
noncomputable def gradingLatGen (k : ℤ) : Submodule ℂ (LatModuleGen L) :=
  ⨆ α : L.carrier, bigradα L (k - halfNorm L α, α)
\end{leancode}

\begin{dictionary}
\lean{L.carrier →₀ ℂ} (finitely supported functions $L\to\C$) is the vector space $\C\{L\}$, and
\lean{eGen L α} is the basis vector $e^\alpha$. \lean{⊗[ℂ]} is the algebraic tensor product over
$\C$ and \lean{p ⊗ₜ x} an elementary tensor. \lean{latChargeGen L α} is $1\otimes e^\alpha$.
\lean{mkEα L α} is the map $p\mapsto p\otimes e^\alpha$, and \lean{bigradα L (w,α)} is its image of
$(F_r)_w$, i.e.\ $(F_r)_w\otimes e^\alpha$. \lean{gradingLatGen L k} is the sum (\lean{⨆}, the
supremum of subspaces) of these over all $\alpha$ with $w=k-\tfrac12\langle\alpha,\alpha\rangle$. That
this sum is direct and that the spaces exhaust $V_L$ is proved in Chapter~\ref{ch:03}.
\end{dictionary}

\begin{physics}
$V_L=F_r\otimes\C\{L\}$ is the full chiral state space of the compact boson on $\mathbb R^r/L$: oscillators tensored with a momentum eigenstate $e^\alpha=\,:\!e^{i\alpha\cdot X}\!:(0)\,|0\rangle$. Under the state--operator correspondence, $p\otimes e^\alpha$ is the local operator obtained by dressing the vertex operator $:\!e^{i\alpha\cdot X}\!:$ with derivatives of $X$; its $L_0$ eigenvalue is the oscillator level plus the momentum energy $\tfrac12\langle\alpha,\alpha\rangle=\tfrac12 p^2$. This is a statement about the Hilbert space only; the OPEs are defined in Chapter~\ref{ch:03}.
\end{physics}

\subsection{Formal distributions and modes}

\begin{definition}[Formal distributions, modes, $n$-th products]\label{ch04:def:mode}
A \emph{formal distribution} with values in a vector space $\mathcal W$ is a family
$f=(f_t)_{t\in\Z}$, thought of as $f(z)=\sum_t f_t z^t$. For a field
$Y(u,z)=\sum_n u_{(n)}z^{-n-1}$ the \emph{$n$-th mode} $u_{(n)}$ is the coefficient of
$z^{-n-1}$, and the \emph{$n$-th product} of a vertex algebra is $u_{(n)}v$.
\end{definition}

In Lean these are \lean{FormalDist 𝒲} $=$ \lean{ℤ → 𝒲} (\leanpath{VOA/Foundations/FormalSeries.lean}),
\lean{mode f m = f (-m - 1)} (\leanpath{VOA/Foundations/FormalCalculus.lean}),
\lean{V.smul u n v} $:=$ \lean{mode (V.Y u) n v} (\leanpath{VOA/Core/VertexAlgebra.lean}) and
\lean{W.L n} $:=$ \lean{mode (W.Y W.conformalVector) (n + 1)}
(\leanpath{VOA/Core/VertexOperatorAlgebra.lean}); they are quoted verbatim in
Definitions~\ref{l01:def:formaldist}, \ref{l01:def:mode}, \ref{l03:def:smulT} and
Section~\ref{l04:sec:conv}.

\begin{dictionary}
\lean{f t} is the coefficient of $z^t$; hence \lean{mode f n = f (-n-1)} is the coefficient of
$z^{-n-1}$, i.e.\ \lean{mode (Y u) n} is $u_{(n)}$. \lean{V.smul u n v} is $u_{(n)}v$ (the name
``smul'' is historical and has nothing to do with scalar multiplication). \lean{W.L n} is
$L_n=\omega_{(n+1)}$. The structures \lean{PreVertexAlgebra} and \lean{VertexOperatorAlgebra}
(fields \lean{Y}, \lean{vacuum}, \lean{conformalVector}, \lean{grading}, \lean{centralCharge} and
the axioms) are defined in Definitions~\ref{l03:def:VA} and~\ref{l03:def:VOA}; \lean{W.grading k} is the weight space $W_k$, a
subspace of $W$.
\end{dictionary}

\begin{physics}
A formal distribution $Y(u,z)=\sum u_{(n)}z^{-n-1}$ is a chiral quantum field expanded in its Laurent modes, and $u_{(n)}v$ is the coefficient of $(z-w)^{-n-1}$ in the OPE $Y(u,z)v(w)$. The negative products ($n<0$) are the regular (normal-ordered) terms, and $n\ge0$ the singular ones, so $u_{(0)}v$ is the simple-pole term, $u_{(1)}v$ the double pole, and so on. Identity $L_n=\omega_{(n+1)}$ says $L_n$ are the Laurent modes of the stress tensor $T(z)=\sum L_nz^{-n-2}$.
\end{physics}

\subsection{The vertex operators of charged vacua}

The heart of the lattice construction is the vertex operator of $e^\alpha$. Embed $L$ in $\mathbb{R}^r$
by an orthonormal frame, so that each $\alpha$ has real coordinates $\alpha_b$ ($b=0,\dots,r-1$)
with $\sum_b\alpha_b\beta_b=\langle\alpha,\beta\rangle$. Put $\alpha(n)=\sum_b\alpha_b h_b(n)$. Then
\begin{equation}\label{ch04:eq:Yealpha}
  Y(e^\alpha,z)\,(p\otimes e^\beta)
  =\varepsilon(\alpha,\beta)\, z^{\langle\alpha,\beta\rangle}\,
  E^-(\alpha,z)E^+(\alpha,z)\,p\otimes e^{\alpha+\beta},
\end{equation}
where
\begin{gather*}
  E^-(\alpha,z)=\exp\Big(\sum_{n\ge1}\frac{\alpha(-n)}{n}z^n\Big)=\sum_{N\ge0}S_N(\alpha)z^N,
  \\
  E^+(\alpha,z)=\exp\Big(-\sum_{n\ge1}\frac{\alpha(n)}{n}z^{-n}\Big)=\sum_{N\ge0}E^+(\alpha)_{-N}z^{-N},
\end{gather*}
the $S_N(\alpha)\in F_r$ are the \emph{Schur polynomials} in the $\alpha(-n)$ (for example
$S_1(\alpha)=\alpha(-1)$, $S_2(\alpha)=\tfrac12(\alpha(-1)^2+\alpha(-2))$), and $\varepsilon$ is the
cocycle of \S\ref{ch04:sec:cocycle}. The product $EE(\alpha,z):=E^-(\alpha,z)E^+(\alpha,z)$ is the
``boson dressing''. Physically, \eqref{ch04:eq:Yealpha} is the normal-ordered exponential
$:\!e^{i\alpha\cdot X(z)}\!:$ of the free bosons times the momentum-shift operator $e^\alpha
z^{\alpha(0)}$.

\begin{leancode}
-- VOA/Examples/LatticeGenChargeFields.lean
noncomputable def chargeCoord (α : L.carrier) (b : Fin (rank L)) : ℂ :=
  (orthCoord L α (rankEquiv L b) : ℂ)
noncomputable def vVar (α : L.carrier) (m : ℕ) : HeisenbergN.FockN (rank L) :=
  ∑ a : Fin (rank L), chargeCoord L α a • (MvPolynomial.X ((m, a)) : HeisenbergN.FockN (rank L))
noncomputable def vSchur (α : L.carrier) : ℕ → HeisenbergN.FockN (rank L)
  | 0 => 1
  | (N+1) => ((N+1 : ℂ))⁻¹ • ∑ m ∈ Finset.range (N+1), vVar L α m * vSchur α (N - m)
noncomputable def aExpVec (α : L.carrier) : ℕ → End ℂ (HeisenbergN.FockN (rank L))
  | 0 => LinearMap.id
  | (N+1) => ((N+1 : ℂ))⁻¹ • ∑ m ∈ Finset.range (N+1),
      (-(chargeOsc L α (m+1 : ℤ))).comp (aExpVec α (N - m))

-- VOA/Examples/LatticeGenVertex.lean
noncomputable def EminusVecField (α : L.carrier) : FormalDist (End ℂ (FockN (rank L))) :=
  fun t => if 0 ≤ t then EminusVec L α t.toNat else 0
noncomputable def EplusVecField (α : L.carrier) : FormalDist (End ℂ (FockN (rank L))) :=
  fun t => if t ≤ 0 then aExpVec L α (-t).toNat else 0
noncomputable def EEVec (α : L.carrier) : FormalDist (End ℂ (FockN (rank L))) :=
  nprod (EminusVecField_isField L α) (EplusVecField_isField L α)

theorem Yeαvec_tmul (α β : L.carrier) (t : ℤ) (p : FockN (rank L)) :
    Yeαvec L α t (latStateGen L p β)
      = latticeCocycle L α β • latStateGen L (EEVec L α (t - L.form α β) p) (α + β)
\end{leancode}

\begin{dictionary}
\lean{chargeCoord L α b} is the real coordinate $\alpha_b$ (cast to $\C$);
\lean{vVar L α m} $=\sum_b\alpha_b x_{(m,b)}=\alpha(-m-1)\chfourone$; \lean{vSchur L α N} is $S_N(\alpha)$,
defined by Newton's recursion $(N+1)S_{N+1}=\sum_{m=0}^{N}\alpha(-m-1)S_{N-m}$, which is the
coefficientwise form of $\tfrac{d}{dz}E^-=(\sum_n\alpha(-n)z^{n-1})E^-$. \lean{EminusVec L α N} is
multiplication by $S_N(\alpha)$, and \lean{aExpVec L α N} is $E^+(\alpha)_{-N}$ (same recursion with
$-\alpha(m+1)$; \lean{chargeOsc L α m} is the operator $\alpha(m)$). Thus
\lean{EminusVecField L α t} is the coefficient of $z^t$ in $E^-(\alpha,z)$ (zero for $t<0$) and
\lean{EplusVecField L α t} the coefficient of $z^t$ in $E^+(\alpha,z)$ (zero for $t>0$).
\lean{nprod} is the normal-ordered product $:\!a(z)b(z)\!:$ of two fields in the same variable,
defined coefficientwise (Chapter~\ref{ch:01}); here, with $a=E^-$ (only creation operators) and
$b=E^+$ (only annihilation operators), it is simply the product $E^-(\alpha,z)E^+(\alpha,z)$.
Finally \lean{Yeαvec L α t} is the coefficient of $z^t$ in $Y(e^\alpha,z)$, and \lean{Yeαvec_tmul}
says it is $\varepsilon(\alpha,\beta)$ times the coefficient of $z^{t-\langle\alpha,\beta\rangle}$ of
$EE(\alpha,z)$: this is \eqref{ch04:eq:Yealpha} read coefficientwise.
\end{dictionary}

The full vertex map \lean{YfullGen L hEven hPos} of $V_L$ is produced in Chapter~\ref{ch:03} by a
reconstruction theorem. For this chapter we need exactly one fact about it: on spanning ``word
states'' it is given by explicit fields. A \emph{word} is a finite list
$l=[(i_1,a_1),\dots,(i_k,a_k)]$, of length $|l|=k$, and the \emph{word state} is
\begin{equation*}
  w(l,\alpha)=h_{a_1}(-i_1-1)\cdots h_{a_k}(-i_k-1)\,e^\alpha .
\end{equation*}
The word states span $V_L$ (\lean{wordStateLatGen_span}). The \emph{word field} is the iterated
normal-ordered product
\begin{equation*}
  Y_{\mathrm{word}}(l,\alpha)(z)
  =\;:\!\tfrac{1}{i_1!}\partial^{i_1}h_{a_1}(z)\;:\!\cdots
  :\!\tfrac{1}{i_k!}\partial^{i_k}h_{a_k}(z)\,Y(e^\alpha,z)\!:\cdots\!:\;,
\end{equation*}
where $h_a(z)=\sum_n h_a(n)z^{-n-1}$ is the current on $V_L$ \emph{including} its zero mode
$h_a(0)$, which acts on $p\otimes e^\beta$ by the momentum coordinate $\beta_a$.

\begin{leancode}
-- VOA/Examples/LatticeGenVOAInstance.lean
theorem YfullGen_word (hEven : L.Even) (hPos : L.PosDef) (l : List (ℕ × Fin (rank L)))
    (α : L.carrier) :
    YfullGen L hEven hPos (wordStateLatGen L l α) = YwordLatGen L l α
\end{leancode}

The lattice VOA is \lean{latticeVOAGen L hEven hPos} (Chapter~\ref{ch:03}), with
\lean{grading := gradingLatGen L}, \lean{Y := YfullGen L hEven hPos}, vacuum $e^0$ and conformal
vector $\omega=\tfrac12\sum_a h_a(-1)^2\chfourone\otimes e^0$ (\lean{omegaLatGen L}); and
\begin{leancode}
-- VOA/Moonshine/LeechLatticeVOA.lean
noncomputable def leechVOA : VertexOperatorAlgebra (LatGen.LatModuleGen leech) :=
  LatGen.latticeVOAGen leech leech_even leech_posDef
theorem leechVOA_centralCharge : leechVOA.centralCharge = 24
theorem leechVOA_grading_one :
    leechVOA.grading 1 = LatGen.bigradα leech (1, 0)
\end{leancode}
The last statement, proved in Chapter~\ref{ch:03} from \lean{leech_no_roots}, says
$(V_\Leech)_1=(F_{24})_1\otimes e^0=\operatorname{span}\{h_a(-1)\chfourone\otimes e^0\}$: since
$\Leech$ has no norm-$2$ vectors, no $e^\alpha$ has weight $1$.

\subsection{The cocycle}\label{ch04:sec:cocycle}

\begin{definition}[The basis cocycle]\label{ch04:def:cocycle}
Fix a $\Z$-basis $B_1,\dots,B_r$ of $L$ and write $\alpha=\sum_i\alpha_iB_i$ with $\alpha_i\in\Z$.
The \emph{cocycle exponent} and \emph{cocycle} are
\begin{equation*}
  c(\alpha,\beta)=\sum_{j<i}\langle B_i,B_j\rangle\,\alpha_i\beta_j\in\Z,\qquad
  \varepsilon(\alpha,\beta)=(-1)^{c(\alpha,\beta)}\in\{\pm1\}.
\end{equation*}
\end{definition}

\begin{leancode}
-- VOA/Foundations/LatticeCocycle.lean
noncomputable def cocExp (α β : L.carrier) : ℤ :=
  ∑ i : Fin n, ∑ j : Fin n,
    (if j < i then L.form (B i) (B j) else 0) * B.repr α i * B.repr β j
noncomputable def cocycle (α β : L.carrier) : ℂ := (-1 : ℂ) ^ (cocExp L B α β)
noncomputable def latticeBasis :
    Basis (Fin (Fintype.card (Module.Free.ChooseBasisIndex ℤ L.carrier))) ℤ L.carrier :=
  (Module.Free.chooseBasis ℤ L.carrier).reindex (Fintype.equivFin _)
noncomputable def latticeCocycle (α β : L.carrier) : ℂ := cocycle L (latticeBasis L) α β

theorem latticeCocycle_add_left (α α' β : L.carrier) :
    latticeCocycle L (α + α') β = latticeCocycle L α β * latticeCocycle L α' β
theorem latticeCocycle_add_right (α β β' : L.carrier) :
    latticeCocycle L α (β + β') = latticeCocycle L α β * latticeCocycle L α β'
theorem latticeCocycle_sq (α β : L.carrier) :
    latticeCocycle L α β * latticeCocycle L α β = 1
theorem latticeCocycle_comm (hEven : L.Even) (α β : L.carrier) :
    latticeCocycle L α β * latticeCocycle L β α = (-1 : ℂ) ^ (L.form α β)
\end{leancode}

\begin{dictionary}
\lean{B.repr α i} is the $i$-th coordinate $\alpha_i$ of $\alpha$ in the basis $B$.
\lean{latticeBasis L} is \emph{some} basis, obtained from Mathlib's \lean{Module.Free.chooseBasis}
--- a choice made by the axiom of choice, which cannot be computed with. Every property of
$\varepsilon$ the library uses is proved for an \emph{arbitrary} basis, so the choice is harmless
for correctness; but a quantity depending on the basis itself (we meet one in
\S\ref{ch04:sec:form}, namely $\varepsilon(\alpha,\alpha)$) cannot be evaluated. $\varepsilon$ is
bimultiplicative (being $(-1)$ to a bilinear form), a sign, and for even $L$ satisfies the
commutator relation $\varepsilon(\alpha,\beta)\varepsilon(\beta,\alpha)=(-1)^{\langle\alpha,\beta\rangle}$,
because $c(\alpha,\beta)+c(\beta,\alpha)=\langle\alpha,\beta\rangle-\sum_i\langle B_i,B_i\rangle
\alpha_i\beta_i$ and the diagonal correction is even.
\end{dictionary}

\begin{physics}
The sign $\varepsilon(\alpha,\beta)$ is the standard Klein-factor (cocycle) correction in the bosonic construction of lattice vertex operators: without it, vertex operators $:\!e^{i\alpha\cdot X}\!:$ and $:\!e^{i\beta\cdot X}\!:$ would commute rather than anticommute when $\langle\alpha,\beta\rangle$ is odd, violating locality. Physically the cocycle is a gauge-like convention, and different choices of basis give isomorphic theories; this is why only the commutator $(-1)^{\langle\alpha,\beta\rangle}$ matters. Heuristic.
\end{physics}

\begin{example}[Two bases, two cocycles, one commutator]\label{ch04:ex:cocycle}
Take the $A_2$ root lattice, with basis $B_1,B_2$ and Gram matrix
$\left(\begin{smallmatrix}2&-1\\-1&2\end{smallmatrix}\right)$. The only off-diagonal entry with
$j<i$ is $\langle B_2,B_1\rangle=-1$, so $c(\alpha,\beta)=-\alpha_2\beta_1$ and
$\varepsilon(\alpha,\beta)=(-1)^{\alpha_2\beta_1}$. Thus $\varepsilon(B_1,B_2)=1$ and
$\varepsilon(B_2,B_1)=-1$. If instead the basis is listed in the opposite order ($B_1'=B_2$,
$B_2'=B_1$), the same recipe gives $\varepsilon'(\alpha,\beta)=(-1)^{\alpha_1\beta_2}$, so
$\varepsilon'(B_1,B_2)=-1$ and $\varepsilon'(B_2,B_1)=1$: a \emph{different function}. But the
commutator is the same, $\varepsilon(B_1,B_2)\varepsilon(B_2,B_1)=\varepsilon'(B_1,B_2)\varepsilon'(B_2,B_1)
=-1=(-1)^{\langle B_1,B_2\rangle}$, and so are bimultiplicativity, $\varepsilon^2=1$ and
$\varepsilon(-\alpha,\beta)=\varepsilon(\alpha,\beta)$. These are the only properties of $\varepsilon$
that the lattice VOA and this chapter use, and they hold for every basis. This is what ``the choice
is harmless'' means: changing the basis changes $\varepsilon$, and the resulting VOAs are isomorphic,
but no theorem quoted here depends on which basis was chosen. A quantity like
$\varepsilon(\alpha,\alpha)$ does depend on the basis in general, and \S\ref{ch04:sec:form} is careful
never to assume a value for it.
\end{example}

%======================================================================
\section{The Fock involution \texorpdfstring{$\sigma$}{sigma}}\label{ch04:sec:sigma}

The reflection $X\mapsto-X$ negates every oscillator. On the Fock space it is therefore the
algebra automorphism negating every variable.

\begin{definition}[Fock involution]\label{ch04:def:sigma}
$\sigma:F_r\to F_r$ is the unique $\C$-algebra homomorphism with $\sigma(x_{(m,a)})=-x_{(m,a)}$.
\end{definition}

\begin{physics}
$\sigma$ is the action of the reflection $X^a\mapsto-X^a$ on the oscillator part of the Hilbert space: each creation operator $\alpha^a_{-n}$ flips sign, so a state with $k$ quanta acquires $(-1)^k$. This is the oscillator half of the $\mathbb Z_2$ global symmetry of the free bosons, and it is the reflection that will be gauged in the orbifold.
\end{physics}

\begin{proposition}\label{ch04:prop:sigma}
(a) $\sigma^2=1$. (b) $\sigma\,h_a(n)=-h_a(n)\,\sigma$ for all $a$ and all $n\in\Z$.
(c) On a monomial, $\sigma(x^d)=(-1)^{\deg d}x^d$.
\end{proposition}

\begin{leancode}
-- VOA/Moonshine/LatticeInvolution.lean
noncomputable def sigmaAlg : FockN r →ₐ[ℂ] FockN r := MvPolynomial.aeval (fun v => -(X v : FockN r))

theorem sigmaAlg_involutive : Function.Involutive (sigmaAlg r)

noncomputable def sigmaN : FockN r ≃ₗ[ℂ] FockN r :=
  LinearEquiv.ofInvolutive (sigmaAlg r).toLinearMap (sigmaAlg_involutive r)

theorem sigmaN_oscN (k : ℂ) (a : Fin r) (n : ℤ) (p : FockN r) :
    sigmaN r (oscN r k a n p) = -oscN r k a n (sigmaN r p)

-- VOA/Moonshine/PlusWeightOne.lean
theorem sigmaN_monomial (d : (ℕ × Fin r) →₀ ℕ) :
    sigmaN r (monomial d 1) = (-1 : ℂ) ^ d.degree • monomial d 1
\end{leancode}

\begin{dictionary}
\lean{→ₐ[ℂ]} is a $\C$-algebra homomorphism; \lean{MvPolynomial.aeval f} is the homomorphism
sending each variable $v$ to $f(v)$, here $-x_v$. \lean{Function.Involutive g} means $g(g(p))=p$ for
all $p$. \lean{≃ₗ[ℂ]} is a $\C$-linear bijection (linear equivalence), and
\lean{LinearEquiv.ofInvolutive} packages an involutive linear map as one, with itself as inverse.
\lean{sigmaN_oscN} is (b), for every level $k$ (we only use $k=1$). \lean{monomial d 1} is $x^d$
and \lean{sigmaN_monomial} is (c). Same as the textbook statements; no hypotheses.
\end{dictionary}

\begin{physics}
Part (b) says the current $i\partial X^a$ is odd under the symmetry, which is exactly how $\sigma$ implements $X\mapsto-X$ on the modes. Part (c) says the symmetry charge of a state is the parity of its number of oscillators; it is the statement that the $\mathbb Z_2$ quantum number is computed by counting quanta.
\end{physics}

\paragraph{Proof.} (a) $\sigma^2$ is an algebra homomorphism fixing every variable, hence the
identity. (b) If $n<0$, $h_a(n)$ is multiplication by a variable $x$, and
$\sigma(xp)=\sigma(x)\sigma(p)=-x\,\sigma(p)$. If $n>0$, $h_a(n)$ is a multiple of a partial
derivative $\partial_v$, and $\sigma\partial_v=-\partial_v\sigma$: both $\sigma\partial_v\sigma^{-1}$
and $-\partial_v$ are derivations of $F_r$, and they agree on the generators
($\sigma\partial_v\sigma^{-1}x_w=\sigma\partial_v(-x_w)=-\delta_{vw}$). If $n=0$ both sides vanish.
(c) is multiplicativity: $\sigma(x^d)=\prod(-x_{(m,a)})^{d_{(m,a)}}$.\qed

\begin{status}
Proved, no hypotheses (\leanpath{VOA/Moonshine/LatticeInvolution.lean},
\leanpath{VOA/Moonshine/PlusWeightOne.lean}).
\end{status}

%======================================================================
\section{Lifting \texorpdfstring{$-1$}{-1}: the cocycle needs no correction}\label{ch04:sec:theta}

On $V_L=F_r\otimes\C\{L\}$ the reflection should act by $\sigma$ on oscillators and by
$e^\alpha\mapsto e^{-\alpha}$ on momenta. Before defining this map we should ask whether it has a
chance of commuting with the vertex operators, and \eqref{ch04:eq:Yealpha} shows where the danger
lies: conjugating $Y(e^\alpha,z)$ by such a map produces $Y(e^{-\alpha},z)$ only \emph{up to the
ratio} $\varepsilon(\alpha,\beta)/\varepsilon(-\alpha,-\beta)$ on the sector $e^{\beta}$. For a general
$2$-cocycle this ratio is not $1$, and the textbook theory (Frenkel--Lepowsky--Meurman, \S10.4)
corrects the lift by a sign function $\eta:L\to\{\pm1\}$, $e^\alpha\mapsto\eta(\alpha)e^{-\alpha}$. For
the library's cocycle no correction is needed.

\begin{proposition}\label{ch04:prop:cocneg}
For every integral lattice $L$ and all $\alpha,\beta\in L$:
$\varepsilon(-\alpha,\beta)=\varepsilon(\alpha,\beta)=\varepsilon(\alpha,-\beta)$; hence
$\varepsilon(-\alpha,-\beta)=\varepsilon(\alpha,\beta)$.
\end{proposition}

\begin{leancode}
-- VOA/Moonshine/LatticeInvolution.lean
theorem latticeCocycle_neg_left (α β : L.carrier) :
    latticeCocycle L (-α) β = latticeCocycle L α β
theorem latticeCocycle_neg_right (α β : L.carrier) :
    latticeCocycle L α (-β) = latticeCocycle L α β
theorem latticeCocycle_neg_neg (α β : L.carrier) :
    latticeCocycle L (-α) (-β) = latticeCocycle L α β
\end{leancode}

\begin{dictionary}
No hypothesis on $L$, not even evenness: the statement is about the specific cocycle
\lean{latticeCocycle} of Definition~\ref{ch04:def:cocycle}, and it is literally the textbook
statement ``$\varepsilon$ is invariant under $-1$''. It \emph{would be false} for a general
normalized $2$-cocycle with the same commutator map; that is why the textbook needs $\eta$. The
library escapes because its $\varepsilon$ is \emph{bimultiplicative}, not merely a cocycle.
\end{dictionary}

\begin{physics}
This is a statement about the sign (Klein-factor) conventions only: reflecting both momenta leaves the cocycle sign unchanged. Physically it means the reflection $\alpha\mapsto-\alpha$ does not need an extra momentum-dependent phase to be a symmetry of the operator algebra, which for a general cocycle choice would be required (the correcting function of FLM). A convention check, not a new physical effect.
\end{physics}

\paragraph{Proof.} By multiplicativity in the first slot,
$\varepsilon(\alpha,\beta)\varepsilon(-\alpha,\beta)=\varepsilon(0,\beta)=1$, and
$\varepsilon(\alpha,\beta)^2=1$; so $\varepsilon(-\alpha,\beta)=\varepsilon(\alpha,\beta)$. Equivalently,
$c(-\alpha,\beta)=-c(\alpha,\beta)$ and $(-1)^{-c}=(-1)^c$. The right slot is identical.\qed

\begin{warning}
This is a place where a convention could silently break the construction. If the lattice VOA had
been built with a cocycle that is only a cocycle --- for instance
$\varepsilon'(\alpha,\beta)=\varepsilon(\alpha,\beta)\,\eta(\alpha+\beta)\eta(\alpha)^{-1}\eta(\beta)^{-1}$
for a sign function $\eta$ that is not a character, which gives an isomorphic VOA but is not
bimultiplicative --- then the naive lift below would \emph{not} be an automorphism and one would
need the correcting function of FLM~\S10.4. The project roadmap records this as a known trap
(``a non-bimultiplicative cocycle would need a correcting function on $\Lambda$''). The safeguard
here is that Theorem~\ref{ch04:thm:thetamode} below is proved: whatever one believes about
cocycles, the machine has checked that the map commutes with all vertex operators.
\end{warning}

\begin{status}
Proved, no hypotheses (\leanpath{VOA/Moonshine/LatticeInvolution.lean}).
\end{status}

\begin{definition}[The lift $\theta$ of $-1$]\label{ch04:def:theta}
$\theta:V_L\to V_L$ is the linear map $\theta(p\otimes e^\alpha)=\sigma(p)\otimes e^{-\alpha}$,
i.e.\ $\theta=\sigma\otimes\iota$ where $\iota(e^\alpha)=e^{-\alpha}$.
\end{definition}

\begin{physics}
$\theta$ is the full reflection $X\to-X$ acting on states: oscillators flip sign (this is $\sigma$) and a state of momentum $\alpha$ is sent to momentum $-\alpha$. It is the unitary-type operator implementing a $\mathbb Z_2$ global symmetry of the compactified bosons, and is the lift of the lattice isometry $-1$ to the chiral algebra. Heuristic: on the torus $\mathbb R^{24}/\Leech$ the point-group element $-1$ is a symmetry for every lattice.
\end{physics}

\begin{proposition}\label{ch04:prop:thetabasic}
$\theta^2=1$; $\theta$ fixes the vacuum $e^0$ and the conformal vector $\omega$; and $\theta$
preserves each weight space $(V_L)_k$.
\end{proposition}

\begin{leancode}
-- VOA/Moonshine/LatticeInvolution.lean
noncomputable def thetaLat : LatModuleGen L ≃ₗ[ℂ] LatModuleGen L :=
  TensorProduct.congr (sigmaN (rank L)) (Finsupp.domLCongr (Equiv.neg L.carrier))

theorem thetaLat_tmul (p : FockN (rank L)) (α : L.carrier) :
    thetaLat L (p ⊗ₜ[ℂ] eGen L α) = sigmaN (rank L) p ⊗ₜ[ℂ] eGen L (-α)

theorem thetaLat_involutive : Function.Involutive (thetaLat L)
theorem thetaLat_vacuum : thetaLat L (latVacuumGen L) = latVacuumGen L
theorem thetaLat_omega : thetaLat L (omegaLatGen L) = omegaLatGen L
theorem thetaLat_mem_gradingLatGen {k : ℤ} {x : LatModuleGen L} (hx : x ∈ gradingLatGen L k) :
    thetaLat L x ∈ gradingLatGen L k
\end{leancode}

\begin{dictionary}
\lean{TensorProduct.congr f g} is $f\otimes g$ for linear equivalences $f,g$;
\lean{Finsupp.domLCongr (Equiv.neg L.carrier)} is the linear equivalence of $\C\{L\}$ induced by the
bijection $\alpha\mapsto-\alpha$ of basis indices, i.e.\ $e^\alpha\mapsto e^{-\alpha}$. This is spelled
out by \lean{thetaLat_tmul}, which is Definition~\ref{ch04:def:theta} verbatim, with \emph{no}
correcting sign. \lean{latVacuumGen L} is $e^0$ and \lean{omegaLatGen L} is $\omega$. The four
remaining statements are the four parts of the proposition.
\end{dictionary}

\begin{physics}
These three facts are what one expects of a symmetry in a CFT: it squares to one (a $\mathbb Z_2$), leaves the vacuum invariant (the vacuum is not spontaneously broken), leaves the stress tensor $T=\omega$ invariant (so the symmetry commutes with the Virasoro algebra and with $L_0$), and hence acts on each energy level. The last point is what permits decomposing every weight space into $\pm1$ eigenspaces.
\end{physics}

\paragraph{Proof.} $\theta^2=\sigma^2\otimes\iota^2=1$. $\theta(1\otimes e^0)=\sigma(1)\otimes
e^0=1\otimes e^0$. The conformal vector $\tfrac12\sum_a h_a(-1)^2\chfourone\otimes e^0$ is quadratic
in the oscillators, so $\sigma$ fixes it. Finally $\sigma$ preserves Fock weight and
$\tfrac12\langle-\alpha,-\alpha\rangle=\tfrac12\langle\alpha,\alpha\rangle$, so $(F_r)_w\otimes e^\alpha$
goes to $(F_r)_w\otimes e^{-\alpha}$, of the same conformal weight.\qed

\begin{status}
Proved, no hypotheses (\leanpath{VOA/Moonshine/LatticeInvolution.lean}).
\end{status}
%======================================================================
\section{\texorpdfstring{$\theta$}{theta} commutes with every vertex operator}\label{ch04:sec:intertw}

The central fact about $\theta$ is that it is compatible with the whole vertex algebra structure:
\begin{equation}\label{ch04:eq:thetaY}
  \theta\,Y(u,z)\,\theta^{-1}=Y(\theta u,z)\qquad\text{for all }u\in V_L .
\end{equation}
To state the intermediate steps conveniently we introduce a small piece of language.

\begin{definition}[Intertwining of formal distributions]\label{ch04:def:intertw}
Let $g$ be a linear endomorphism of a vector space $V$, $c\in\C$, and let $A(z),B(z)$ be formal
distributions with coefficients in $\End(V)$. We say that \emph{$g$ intertwines $A$ with $B$ up
to $c$}, written $gA=c\,Bg$, if $g\,A_t\,x=c\,B_t\,g\,x$ for every coefficient index $t$ and every
$x\in V$; that is, $gA(z)=c\,B(z)g$ as formal series.
\end{definition}

\begin{leancode}
-- VOA/Moonshine/LatticeInvolution.lean
def Intertw (g : End ℂ V) (c : ℂ) (A B : FormalDist (End ℂ V)) : Prop :=
  ∀ (t : ℤ) (x : V), g (A t x) = c • B t (g x)
\end{leancode}

\begin{dictionary}
\lean{End ℂ V} is the space of $\C$-linear maps $V\to V$. \lean{A t} is the coefficient of $z^t$
of $A(z)$; the definition is word for word Definition~\ref{ch04:def:intertw}. It is a
coefficientwise identity, so it involves no convergence questions.
\end{dictionary}

\begin{physics}
This is only bookkeeping language: ``$g$ intertwines $A$ with $B$'' says that conjugating the field $A(z)$ by the symmetry operator $g$ gives the field $c\,B(z)$, i.e.\ the field transforms covariantly under the symmetry. It is the mode-level form of the statement that a global symmetry $g$ maps local operators to local operators, $g\,\mathcal O(z)\,g^{-1}=(g\mathcal O)(z)$.
\end{physics}

The calculus of this relation consists of three closure properties, each immediate from the
definition:
\begin{enumerate}[label=(\alph*)]
\item if $gA=c\,Bg$ then $g\,\partial A=c\,(\partial B)\,g$ and $g(eA)=c\,(eB)g$ for scalars $e$;
\item if $ga=c\,a'g$ and $gb=d\,b'g$ for fields $a,b,a',b'$, then
  $g\,{:}ab{:}=cd\,{:}a'b'{:}\,g$ for the normal-ordered products.
\end{enumerate}
For (b): the coefficient of the normal-ordered product is a (locally finite) sum of products
$a_mb_n$ or $b_na_m$ with the same ordering rule on both sides, and $g$ passes through each factor
picking up $c$ and $d$.

\begin{theorem}[$\theta$ intertwines the vertex operators]\label{ch04:thm:thetamode}
Let $L$ be even and positive-definite. Then for all $u,v\in V_L$ and $n\in\Z$,
\begin{equation*}
  \theta\big(u_{(n)}v\big)=(\theta u)_{(n)}(\theta v),
\end{equation*}
equivalently \eqref{ch04:eq:thetaY}.
\end{theorem}

\begin{leancode}
-- VOA/Moonshine/LatticeInvolution.lean
theorem thetaLat_YfullGen (u : LatModuleGen L) (t : ℤ) (x : LatModuleGen L) :
    thetaLat L (YfullGen L hEven hPos u t x)
      = YfullGen L hEven hPos (thetaLat L u) t (thetaLat L x)

theorem thetaLat_mode (u : LatModuleGen L) (n : ℤ) (v : LatModuleGen L) :
    thetaLat L (mode (YfullGen L hEven hPos u) n v)
      = mode (YfullGen L hEven hPos (thetaLat L u)) n (thetaLat L v)
\end{leancode}

\begin{dictionary}
Section variables: \lean{(L : IntegralLattice) (hEven : L.Even) (hPos : L.PosDef)}. These are
exactly the hypotheses under which the lattice VOA exists (Chapter~\ref{ch:03}), so they cost
nothing. \lean{YfullGen L hEven hPos u t} is the coefficient of $z^t$ of $Y(u,z)$ for the
\emph{actual} vertex map of the lattice VOA \lean{latticeVOAGen L hEven hPos} (not a model of it),
and \lean{mode (YfullGen … u) n v} is $u_{(n)}v$. The first statement is \eqref{ch04:eq:thetaY}
coefficientwise, the second is the mode form. All variables are universally quantified. The
statement is the textbook one; it could only have been vacuous if the lattice VOA were degenerate,
which it is not (Chapter~\ref{ch:03} computes its graded dimension).
\end{dictionary}

\begin{physics}
This is the statement that $\theta$ is a symmetry of the full OPE algebra: $\theta\,Y(u,z)\,\theta^{-1}=Y(\theta u,z)$ says that applying the reflection to a correlator is the same as reflecting each inserted operator, so all OPE data (all structure constants) are invariant under $X\to-X$. For the vertex operators, $\theta:\;:\!e^{i\alpha\cdot X}\!:\;\mapsto\;:\!e^{-i\alpha\cdot X}\!:$ and $\partial X\mapsto-\partial X$. The point of the theorem is that no extra anomalous phase appears; this holds because the cocycle is bimultiplicative (Proposition~\ref{ch04:prop:cocneg}).
\end{physics}

\paragraph{Proof.} The usual textbook proof goes through Goddard's uniqueness theorem: the field
$\theta Y(u,z)\theta^{-1}$ is local with respect to all $Y(v,z)$ and creates $\theta u$ from the vacuum,
hence equals $Y(\theta u,z)$. We give instead a direct proof, which needs no uniqueness theorem,
because the vertex map of $V_L$ is \emph{known} on a spanning set: $Y(w(l,\alpha),z)$ is the word
field $Y_{\mathrm{word}}(l,\alpha)$. Both sides of \eqref{ch04:eq:thetaY} are linear in $u$, and the
word states span $V_L$, so it suffices to prove, for every word $l$ and every $\alpha$,
\begin{equation}\label{ch04:eq:wordconj}
  \theta\,Y_{\mathrm{word}}(l,\alpha)\,\theta^{-1}=(-1)^{|l|}\,Y_{\mathrm{word}}(l,-\alpha)
  \qquad\text{and}\qquad \theta\,w(l,\alpha)=(-1)^{|l|}w(l,-\alpha),
\end{equation}
since then $\theta Y(w(l,\alpha),z)\theta^{-1}=(-1)^{|l|}Y(w(l,-\alpha),z)=Y(\theta w(l,\alpha),z)$.

\emph{Step 1: the dressing.} $\sigma S_N(\alpha)=S_N(-\alpha)$, since $S_N(\alpha)$ is a polynomial in
the $\alpha(-n)$ and $\sigma\alpha(-n)=-\alpha(-n)=(-\alpha)(-n)$. Likewise
$\sigma\,E^+(\alpha)_{-N}=E^+(-\alpha)_{-N}\,\sigma$, because $E^+(\alpha)_{-N}$ is a polynomial in the
annihilation operators $\alpha(n)$, and $\sigma\alpha(n)=-\alpha(n)\sigma=(-\alpha)(n)\sigma$ by
Proposition~\ref{ch04:prop:sigma}(b). Hence by (b) of the calculus,
$\sigma\,EE(\alpha,z)=EE(-\alpha,z)\,\sigma$.

\emph{Step 2: the charged vacua.} On $p\otimes e^\beta$, using \eqref{ch04:eq:Yealpha},
\begin{align*}
  \theta\,Y(e^\alpha,z)(p\otimes e^\beta)
  &=\varepsilon(\alpha,\beta)z^{\langle\alpha,\beta\rangle}\,\sigma EE(\alpha,z)p\otimes e^{-\alpha-\beta}
  =\varepsilon(\alpha,\beta)z^{\langle\alpha,\beta\rangle}\,EE(-\alpha,z)\sigma p\otimes e^{-\alpha-\beta},\\
  Y(e^{-\alpha},z)\,\theta(p\otimes e^\beta)
  &=\varepsilon(-\alpha,-\beta)z^{\langle-\alpha,-\beta\rangle}\,EE(-\alpha,z)\sigma p\otimes e^{-\alpha-\beta}.
\end{align*}
The two agree by Proposition~\ref{ch04:prop:cocneg} and $\langle-\alpha,-\beta\rangle=\langle\alpha,\beta\rangle$.
So $\theta Y(e^\alpha,z)\theta^{-1}=Y(e^{-\alpha},z)$: the case $l=[\,]$ of \eqref{ch04:eq:wordconj}.

\emph{Step 3: the currents.} $\theta h_a(z)\theta^{-1}=-h_a(z)$, including the zero mode: for
$n\neq0$ this is Proposition~\ref{ch04:prop:sigma}(b), and $h_a(0)$ acts on $p\otimes e^\beta$ by
$\beta_a$, while $\theta$ moves $p\otimes e^\beta$ to the sector $-\beta$, where $h_a(0)$ acts by
$-\beta_a$. By (a), $\theta\big(\tfrac1{i!}\partial^ih_a\big)\theta^{-1}=-\tfrac1{i!}\partial^ih_a$.

\emph{Step 4: induction on the word.} By (b), each additional letter of the word contributes one
factor $-1$, which gives the first half of \eqref{ch04:eq:wordconj}. The second half is
$\theta\big(h_{a_1}(-i_1-1)\cdots e^\alpha\big)=(-1)^k\,h_{a_1}(-i_1-1)\cdots e^{-\alpha}$, by
Proposition~\ref{ch04:prop:sigma}(b) once more.\qed

\paragraph{Remark (what the proof uses).} The argument uses three things: the cocycle identity of
Proposition~\ref{ch04:prop:cocneg}, the explicit form \eqref{ch04:eq:Yealpha} of $Y(e^\alpha,z)$, and the
fact that the vertex map is determined on word states by word fields. It does not use locality or
any uniqueness theorem. Physically, Step 2 is the statement that
$:\!e^{i\alpha\cdot X}\!:\;\mapsto\;:\!e^{-i\alpha\cdot X}\!:$ under $X\mapsto -X$, and Step 3 that
$\partial X\mapsto-\partial X$.

\begin{status}
\lean{thetaLat_mode}: proved, unconditional given \lean{hEven}, \lean{hPos};
\texttt{\#print axioms} gives \texttt{[propext, Classical.choice, Quot.sound]} (checked while
writing this chapter). File \leanpath{VOA/Moonshine/LatticeInvolution.lean}.
\end{status}

%======================================================================
\section{\texorpdfstring{$\theta$}{theta} as an automorphism of order two}\label{ch04:sec:aut}

\begin{definition}[VOA automorphism]\label{ch04:def:aut}
An \emph{automorphism} of a VOA $W$ is a linear bijection $g:W\to W$ with
$g(u_{(n)}v)=(gu)_{(n)}(gv)$ for all $u,v,n$, $g\chfourone=\chfourone$ and $g\omega=\omega$. The
automorphisms form a group $\Aut(W)$ under composition.
\end{definition}

In Lean this is the structure \lean{Aut W} (\leanpath{VOA/Core/Automorphism.lean}), quoted
verbatim and decoded in Definition~\ref{l04:def:aut}: a linear equivalence \lean{toLinearEquiv}
with fields \lean{map_smul'} ($g(u_{(n)}v)=(gu)_{(n)}(gv)$ for all $u,n,v$), \lean{map_vacuum}
and \lean{map_omega}.

\begin{dictionary}
Exactly Definition~\ref{ch04:def:aut} (Gannon \S5.2.1). Fixing $\omega$ forces $g$ to commute with
all $L_n=\omega_{(n+1)}$, hence to preserve the grading. The group structure (\lean{*} is
composition, \lean{1} the identity) is defined in the same file; two automorphisms are equal as
soon as their underlying linear maps are (\lean{Aut.ext}).
\end{dictionary}

\begin{physics}
An automorphism of a VOA is a global (chiral) symmetry of the 2d CFT: a linear map on states preserving the vacuum, the stress tensor and all OPEs, so that it commutes with $L_n$ and acts on correlators by transforming every insertion. The group $\Aut(W)$ is the group of such symmetries; for the Moonshine module it is the Monster, but here we are only constructing one $\mathbb Z_2$ in the abelian-torus theory.
\end{physics}

\begin{theorem}\label{ch04:thm:thetaAut}
Let $L$ be even and positive-definite. Then $\theta\in\Aut(V_L)$, $\theta^2=1$, and if $L\neq0$
then $\theta\neq1$. For the Leech lattice, $\theta_\Leech:=\theta\in\Aut(V_\Leech)$ satisfies
$\theta_\Leech^2=1\neq\theta_\Leech$.
\end{theorem}

\begin{leancode}
-- VOA/Moonshine/LatticeInvolution.lean
noncomputable def thetaAut : Aut (latticeVOAGen L hEven hPos) where
  toLinearEquiv := thetaLat L
  map_smul' u n v := thetaLat_mode L hEven hPos u n v
  map_vacuum := thetaLat_vacuum L
  map_omega := thetaLat_omega L

theorem thetaAut_sq : thetaAut L hEven hPos * thetaAut L hEven hPos = 1
theorem thetaAut_ne_one (hL : ∃ α : L.carrier, α ≠ 0) : thetaAut L hEven hPos ≠ 1

noncomputable def leechTheta : Aut leechVOA := thetaAut leech leech_even leech_posDef
theorem leechTheta_sq : leechTheta * leechTheta = 1 := thetaAut_sq leech leech_even leech_posDef
theorem leechTheta_ne_one : leechTheta ≠ 1
\end{leancode}

\begin{dictionary}
The definition of \lean{thetaAut} is shown in full because it is informative: the four fields are
the linear map of Definition~\ref{ch04:def:theta} and the three proofs
Theorem~\ref{ch04:thm:thetamode} and Proposition~\ref{ch04:prop:thetabasic}. Since
\lean{W.smul u n v} is by definition \lean{mode (W.Y u) n v} and \lean{(latticeVOAGen …).Y} is
\lean{YfullGen …}, the field \lean{map_smul'} is literally \lean{thetaLat_mode}. The hypothesis
\lean{hL} of \lean{thetaAut_ne_one} says $L\neq0$; without it the statement is false ($L=0$ gives
$V_L=F_0=\C$ and $\theta=1$). For $\Leech$ it is supplied by rank $24$. \lean{leechVOA} is
\lean{latticeVOAGen leech leech_even leech_posDef}, so \lean{leechTheta} is literally $\theta$ for
$\Leech$.
\end{dictionary}

\begin{physics}
So $\theta$ is a genuine $\mathbb Z_2$ global symmetry of the torus CFT on $\mathbb R^{24}/\Leech$, of order exactly two (it acts nontrivially, e.g.\ $\partial X^a\mapsto-\partial X^a$, so it is not secretly trivial). This is precisely the symmetry one would orbifold by in the ``$\mathbb R^{24}/(\Leech\rtimes\mathbb Z_2)$'' construction of the Moonshine module.
\end{physics}

\paragraph{Proof.} The automorphism property is Theorem~\ref{ch04:thm:thetamode} together with
Proposition~\ref{ch04:prop:thetabasic}; $\theta^2=1$ is the involutivity of
Proposition~\ref{ch04:prop:thetabasic}. If $\alpha\neq0$ then $\theta e^\alpha=e^{-\alpha}$, and
$e^{-\alpha}\neq e^{\alpha}$ because $-\alpha\neq\alpha$: a lattice is torsion-free, so $2\alpha=0$
forces $\alpha=0$.\qed

\begin{status}
\lean{thetaAut}, \lean{thetaAut_sq}, \lean{thetaAut_ne_one}, \lean{leechTheta},
\lean{leechTheta_sq}, \lean{leechTheta_ne_one}: proved; \texttt{\#print axioms} for
\lean{thetaAut_sq}, \lean{thetaAut_ne_one}, \lean{leechTheta_ne_one} is
\texttt{[propext, Classical.choice, Quot.sound]}. Unconditional (the hypotheses \lean{hEven},
\lean{hPos}, \lean{hL} are discharged for $\Leech$).
\end{status}

%======================================================================
\section{The fixed-point algebra \texorpdfstring{$V_\Lambda^+$}{V+}}\label{ch04:sec:plus}

\begin{definition}[Sub-VOA; fixed-point sub-VOA]\label{ch04:def:orbifold}
A \emph{vertex operator subalgebra} of a VOA $W$ is a subspace $S$ containing $\chfourone$ and
$\omega$, closed under all $n$-th products, and \emph{graded}: if $v\in S$ then each homogeneous
component of $v$ lies in $S$. It is itself a VOA with the restricted vertex map, the same vacuum,
conformal vector and central charge, and grading $S_n=S\cap W_n$. For a subgroup
$G\le\Aut(W)$ the \emph{fixed-point subalgebra} (or \emph{orbifold}) is
$W^G=\{v\in W: gv=v\ \forall g\in G\}$.
\end{definition}

\begin{leancode}
-- VOA/Modularity/Twisted.lean
structure VOASubalgebra (W : VertexOperatorAlgebra 𝒱) where
  carrier : Submodule ℂ 𝒱
  vacuum_mem : W.vacuum ∈ carrier
  conformal_mem : W.conformalVector ∈ carrier
  smul_mem : ∀ u ∈ carrier, ∀ (n : ℤ), ∀ v ∈ carrier, W.smul u n v ∈ carrier
  graded : ∀ (n : ℤ), ∀ v ∈ carrier, piHom W n v ∈ carrier

def gradingS (n : ℤ) : Submodule ℂ S.carrier := (W.grading n).comap S.carrier.subtype

def orbifold (G : Subgroup (Aut W)) : VOASubalgebra W where
  carrier :=
    { carrier := {v | ∀ g ∈ G, g v = v}
      ... }
  ...

noncomputable def orbifoldVOA (G : Subgroup (Aut W)) :
    VertexOperatorAlgebra (orbifold G).carrier :=
  (orbifold G).toVOA

theorem orbifoldVOA_centralCharge (G : Subgroup (Aut W)) :
    (orbifoldVOA G).centralCharge = W.centralCharge := rfl
\end{leancode}

\begin{dictionary}
In the display, ``\lean{...}'' marks omitted \emph{proof} fields of \lean{orbifold} (that the set
is a subspace, contains $\chfourone$ and $\omega$, is closed under products and graded); the only data
is the carrier $\{v\mid\forall g\in G,\ gv=v\}$, exactly $W^G$. \lean{piHom W n v} is the
weight-$n$ component of $v$. \lean{(orbifold G).carrier} is a \lean{Submodule}; as a type,
\lean{↥(orbifold G).carrier} (written with the coercion arrow \lean{↥}) is the set of its elements
$v$ packaged with their membership proof, and the VOA \lean{orbifoldVOA G} lives on that type. Its
weight spaces are \lean{gradingS n}, the elements of $S$ lying in $W_n$ (\lean{comap} along the
inclusion \lean{subtype}). \lean{toVOA} (Definition~\ref{l04:def:subalg}) restricts all structure. The central
charge equality is definitional (\lean{rfl}).
\end{dictionary}

\begin{physics}
Taking $W^G$ is the projection of the Hilbert space onto $G$-invariant states, the first step of orbifolding a CFT by a global symmetry group $G$. The invariant states close under the OPE and contain the identity and stress tensor, so the result is again a chiral algebra with the same central charge $c$; it is a smaller algebra, whose missing operators are recovered (as a different representation) in the twisted sectors. Heuristic: one also has to add the twisted sectors to get a modular invariant theory, which is the content of the later chapters.
\end{physics}

\begin{definition}[$V_\Leech^+$]\label{ch04:def:VPlus}
$V_\Leech^+:=V_\Leech^{\langle\theta\rangle}=\{v\in V_\Leech:\theta v=v\}$, with the VOA structure of
Definition~\ref{ch04:def:orbifold}.
\end{definition}

\begin{leancode}
-- VOA/Moonshine/LatticeInvolution.lean
theorem mem_orbifold_zpowers_iff {𝒱 : Type*} [AddCommGroup 𝒱] [Module ℂ 𝒱]
    {W : VertexOperatorAlgebra 𝒱} (g : Aut W) (v : 𝒱) :
    v ∈ (orbifold (Subgroup.zpowers g)).carrier ↔ g v = v

noncomputable def leechVOAPlus :
    VertexOperatorAlgebra (orbifold (Subgroup.zpowers leechTheta)).carrier :=
  orbifoldVOA (Subgroup.zpowers leechTheta)

theorem mem_leechVOAPlus_iff (v : LatModuleGen leech) :
    v ∈ (orbifold (Subgroup.zpowers leechTheta)).carrier ↔ thetaLat leech v = v :=
  mem_orbifold_zpowers_iff leechTheta v

theorem leechVOAPlus_centralCharge : leechVOAPlus.centralCharge = 24

-- VOA/Moonshine/FrameSectors.lean
abbrev VPlus : Type := ↥(orbifold (Subgroup.zpowers leechTheta)).carrier
\end{leancode}

\begin{dictionary}
\lean{Subgroup.zpowers g} is the cyclic subgroup $\{g^k:k\in\Z\}$ generated by $g$; here
$\{1,\theta\}$. \lean{mem_orbifold_zpowers_iff} says that being fixed by the whole cyclic group is
the same as being fixed by its generator, so \lean{mem_leechVOAPlus_iff} identifies the carrier
with $\{v:\theta v=v\}$, i.e.\ Definition~\ref{ch04:def:VPlus}. \lean{VPlus} is the type of
elements of $V_\Leech^+$ used in later files. The weight spaces of \lean{leechVOAPlus} are
$(V_\Leech^+)_k=V_\Leech^+\cap(V_\Leech)_k$ by \lean{gradingS}. How could this be vacuous? If $\theta$
were the identity, $V_\Leech^+$ would be all of $V_\Leech$ and every subsequent dimension statement
would be wrong; \lean{leechTheta_ne_one} and the dimension $98580\neq196884$ below rule this out.
\end{dictionary}

\begin{physics}
$V_\Leech^+$ is the untwisted $\mathbb Z_2$-invariant sector of the orbifold of the $24$-boson torus theory by the reflection: the operators that survive the projection $P=\tfrac12(1+\theta)$ onto states with even $X\to-X$ parity. Its states are the even combinations such as $\cos(\alpha\cdot X)$ and products of an even number of oscillators; the central charge stays $24$ since the stress tensor $\tfrac12(\partial X)^2$ is even.
\end{physics}

\paragraph{Proof.} If $gv=v$ then $g^kv=v$ for every $k\in\Z$ (for negative $k$ apply $g^{-1}$ to
$gv=v$). The central charge is that of $V_\Leech$, which is $\operatorname{rank}\Leech=24$.\qed

\begin{status}
\lean{leechVOAPlus} is a definition; \lean{mem_leechVOAPlus_iff} and
\lean{leechVOAPlus_centralCharge} are proved, \texttt{\#print axioms}
\texttt{[propext, Classical.choice, Quot.sound]}. Unconditional.
\end{status}

\begin{remark}[Why the fixed points form a VOA, physically and mathematically]
If $\theta u=u$ and $\theta v=v$ then $\theta(u_{(n)}v)=u_{(n)}v$, so $V^+$ is closed under all
products; $\chfourone$ and $\omega$ are fixed; and $\theta$ preserves weights, so components of fixed
vectors are fixed. In CFT language: a symmetry commuting with the stress tensor can be used to
project the theory, and the projected chiral algebra is again a chiral algebra with the same
Virasoro central charge. Nothing more is needed for Definition~\ref{ch04:def:orbifold}. What is
\emph{not} automatic is that $V^+$ is ``large'' --- that it remembers enough of $V$ to be useful
--- and the rest of this chapter makes that precise.
\end{remark}

\begin{remark}[The projection]\label{ch04:rem:proj}
The map $P=\tfrac12(1+\theta):V_\Leech\to V_\Leech^+$ is a projection onto $V_\Leech^+$ (it fixes
$V_\Leech^+$ and kills the $(-1)$-eigenspace $V_\Leech^-$), and $V_\Leech=V_\Leech^+\oplus V_\Leech^-$. On
basis vectors, $P(p\otimes e^\alpha)=\tfrac12(p\otimes e^\alpha+\sigma p\otimes e^{-\alpha})$.
\begin{leancode}
-- VOA/Moonshine/Generation.lean
noncomputable def plusProj : LatModuleGen leech →ₗ[ℂ] VPlus :=
  LinearMap.codRestrict _ ((1 / 2 : ℂ) • (LinearMap.id + (thetaLat leech).toLinearMap))
    (fun v => by simpa using half_add_theta_mem_plus v)
theorem plusProj_of_mem (v : VPlus) : plusProj (v : LatModuleGen leech) = v
\end{leancode}
\lean{LinearMap.codRestrict} regards the map $\tfrac12(1+\theta)$, whose image lies in $V_\Leech^+$,
as a map into $V_\Leech^+$; \lean{plusProj_of_mem} says $P$ is the identity on $V_\Leech^+$
(\lean{(v : LatModuleGen leech)} is $v$ regarded as a vector of $V_\Leech$).
\end{remark}

%======================================================================
\section{Weight one: \texorpdfstring{$(V_\Lambda^+)_1=0$}{(V+)1 = 0}}\label{ch04:sec:wtone}

\begin{lemma}\label{ch04:lem:degwt}
(a) Every monomial satisfies $\deg d\le\operatorname{weight}d$. (b) Consequently $\sigma=-1$ on
$(F_r)_1$.
\end{lemma}

\begin{leancode}
-- VOA/Moonshine/PlusWeightOne.lean
theorem degree_le_weight (d : (ℕ × Fin r) →₀ ℕ) :
    d.degree ≤ Finsupp.weight weightFnN d

theorem sigmaN_weight_one {p : FockN r} (hp : p ∈ gradingWN r 1) : sigmaN r p = -p
\end{leancode}

\begin{dictionary}
\lean{d.degree} and \lean{Finsupp.weight weightFnN d} are the degree and weight of
Definition~\ref{ch04:def:fock}. The hypothesis \lean{hp} says $p\in(F_r)_1$. Same as the textbook
statement.
\end{dictionary}

\begin{physics}
This is the elementary observation that at energy $1$ the only oscillator state is a single quantum $\alpha^a_{-1}|0\rangle$: there is no room for two quanta. A single quantum is odd under $X\to-X$, so the whole level is odd.
\end{physics}

\paragraph{Proof.} (a) Every variable has weight $m+1\ge1$. (b) A weight-$1$ monomial has degree
$\le1$ and is nonconstant, hence has degree $1$; so it is some $x_{(0,a)}=h_a(-1)\chfourone$, and
$\sigma$ multiplies it by $-1$. Extend linearly.\qed

\begin{theorem}[$(V_\Leech^+)_1=0$]\label{ch04:thm:wtone}
The weight-one space of $V_\Leech^+$ is zero.
\end{theorem}

\begin{leancode}
-- VOA/Moonshine/PlusWeightOne.lean
theorem leechVOAPlus_grading_one : leechVOAPlus.grading 1 = ⊥
\end{leancode}

\begin{dictionary}
\lean{leechVOAPlus.grading 1} is $(V_\Leech^+)_1=V_\Leech^+\cap(V_\Leech)_1$ (Definition~\ref{ch04:def:orbifold});
\lean{⊥} is the zero subspace. No hypotheses. This is one half of the statement $\Vnat_1=0$ of the
final theorem; the other half is that the twisted sector starts in weight $3/2$, so that its
$\theta$-fixed part starts in weight $2$ (Chapters~\ref{ch:07}--\ref{ch:08}).
\end{dictionary}

\begin{physics}
There are no weight-one operators left in the orbifold: the only weight-one fields of the torus theory are the $24$ $U(1)$ currents $\partial X^a$ (since $\Leech$ has no norm-$2$ vectors, there are no charged currents $e^{i\alpha\cdot X}$), and they all change sign under the reflection. Hence no continuous symmetry survives the projection, which is the CFT form of the statement that the Monster has no continuous part: $V_1=0$ means no Kac--Moody currents.
\end{physics}

\paragraph{Proof.} By \lean{leechVOA_grading_one}, every $x\in(V_\Leech)_1$ is $p\otimes e^0$ with
$p\in(F_{24})_1$. Then $\theta x=\sigma p\otimes e^0=-x$ by Lemma~\ref{ch04:lem:degwt}. If moreover
$\theta x=x$, then $2x=0$, so $x=0$.\qed

\begin{example}[Weight one, concretely]
$(V_\Leech)_1$ has basis $h_0(-1)e^0,\dots,h_{23}(-1)e^0$, the $24$ ``currents'' of the free bosons;
this is the Lie algebra $\C^{24}$ of the torus $\mathbb{R}^{24}/\Leech$ of continuous symmetries, abelian
since there are no roots. All of them are odd under the reflection. Compare with a lattice
\emph{with} roots, e.g.\ the $E_8$ lattice: there $(V_{E_8})_1$ has dimension $8+240=248$ (the Lie
algebra $E_8$), and the $\theta$-fixed part is spanned by the $120$ vectors $e^\alpha+e^{-\alpha}$
($\alpha$ a root, one per pair $\pm\alpha$), the fixed algebra of the Chevalley involution. The
\emph{absence of roots} is exactly what makes $(V_\Leech^+)_1=0$, and it is the first sign that the
Monster, unlike the symmetry groups of lattice theories, has no continuous part.
\end{example}

\begin{status}
\lean{leechVOAPlus_grading_one}: proved, unconditional, \texttt{\#print axioms}
\texttt{[propext, Classical.choice, Quot.sound]} (\leanpath{VOA/Moonshine/PlusWeightOne.lean}).
\end{status}
%======================================================================
\section{Weight two: 196884 and 98580}\label{ch04:sec:wttwo}

The weight-$2$ space of $\Vnat$ is the Griess algebra, of dimension $196884$. Its untwisted part
is $(V_\Leech^+)_2$. We first compute $(V_\Leech)_2$ and then split it under $\theta$.

\subsection{The weight-two space of a rootless lattice VOA}

A Fock monomial of weight $2$ is either $x_{(1,a)}=h_a(-2)\chfourone$ (one oscillator of weight $2$;
$r$ of them) or $x_{(0,a)}x_{(0,b)}=h_a(-1)h_b(-1)\chfourone$ (two oscillators of weight $1$; one for
each unordered pair $\{a,b\}$ with repetition allowed, so $\binom{r+1}{2}$ of them).

\begin{proposition}\label{ch04:prop:fockwt2}
$\dim(F_r)_2=r+\binom{r+1}{2}$. Moreover $\sigma=-1$ on the linear states $h_a(-2)\chfourone$ and
$\sigma=+1$ on the quadratic states $h_a(-1)h_b(-1)\chfourone$.
\end{proposition}

\begin{leancode}
-- VOA/Moonshine/PlusWeightTwo.lean
theorem finrank_gradingWN_two : finrank ℂ (gradingWN r 2) = r + (r + 1).choose 2

noncomputable def quadFock : Submodule ℂ (FockN r) :=
  Submodule.span ℂ (Set.range (wt2Mon r ∘ Sum.inr))
noncomputable def linFock : Submodule ℂ (FockN r) :=
  Submodule.span ℂ (Set.range (wt2Mon r ∘ Sum.inl))
theorem sigmaN_of_mem_quadFock {p : FockN r} (hp : p ∈ quadFock r) : sigmaN r p = p
theorem sigmaN_of_mem_linFock {p : FockN r} (hp : p ∈ linFock r) : sigmaN r p = -p
\end{leancode}

\begin{dictionary}
\lean{finrank ℂ U} is the dimension of the subspace $U$ over $\C$ (Mathlib's \lean{Module.finrank};
it is $0$ for infinite-dimensional spaces, which is why each finite-dimensionality is proved
separately in the library and why a dimension statement here is a real statement).
\lean{(r + 1).choose 2} is $\binom{r+1}{2}$. \lean{wt2Mon r} enumerates the weight-$2$ monomials by
the disjoint union $\{0..r-1\}\sqcup\mathrm{Sym}^2\{0..r-1\}$: \lean{Sum.inl a} is $h_a(-2)\chfourone$
and \lean{Sum.inr s(a,b)} is $h_a(-1)h_b(-1)\chfourone$; \lean{Submodule.span ℂ (Set.range f)} is the
span of the image of $f$.
\end{dictionary}

\begin{physics}
At level $2$ the oscillator states are of two kinds: one quantum of energy $2$ (the field $\partial^2X^a$) and two quanta of energy $1$ (the bilinears $\partial X^a\partial X^b$, symmetric in $a,b$). The reflection flips the sign of the first and fixes the second, because a state with an even number of quanta is invariant. This sorting is the core of the splitting of the weight-two space into $\theta=\pm1$.
\end{physics}

\begin{theorem}[Weight two of a rootless lattice VOA]\label{ch04:thm:wt2rootless}
Let $L$ be even, positive-definite, with no vectors of norm $2$. Then
\begin{equation*}
  (V_L)_2=(F_r)_2\otimes e^0\;\oplus\bigoplus_{\langle\alpha,\alpha\rangle=4}\C e^\alpha,
  \qquad
  \dim(V_L)_2=r+\binom{r+1}{2}+\#\{\alpha\in L:\langle\alpha,\alpha\rangle=4\}.
\end{equation*}
For the Leech lattice, $\dim(V_\Leech)_2=24+300+196560=196884$.
\end{theorem}

\begin{leancode}
-- VOA/Moonshine/PlusWeightTwo.lean
noncomputable def normFour (hPos : L.PosDef) : Finset L.carrier := (normFour_finite L hPos).toFinset
theorem mem_normFour (hPos : L.PosDef) {α : L.carrier} : α ∈ normFour L hPos ↔ L.form α α = 4
noncomputable def chargeFour (hPos : L.PosDef) : Submodule ℂ (LatModuleGen L) :=
  ⨆ α ∈ normFour L hPos, bigradα L (0, α)

theorem gradingLatGen_two_of_rootless (hEven : L.Even) (hPos : L.PosDef)
    (hR : ∀ v : L.carrier, L.form v v ≠ 2) :
    gradingLatGen L 2 = bigradα L (2, 0) ⊔ chargeFour L hPos

theorem finrank_gradingLatGen_two (hEven : L.Even) (hPos : L.PosDef)
    (hR : ∀ v : L.carrier, L.form v v ≠ 2) :
    finrank ℂ (gradingLatGen L 2)
      = rank L + (rank L + 1).choose 2 + (normFour L hPos).card

theorem leechVOA_finrank_grading_two : finrank ℂ (leechVOA.grading 2) = 196884
\end{leancode}

\begin{dictionary}
\lean{normFour L hPos} is the finite set (\lean{Finset}) of norm-$4$ vectors; finiteness needs
positive-definiteness (\lean{normFour_finite}), and \lean{.card} is its cardinality.
\lean{bigradα L (0, α)} $=\C e^\alpha$, \lean{bigradα L (2, 0)} $=(F_r)_2\otimes e^0$, and \lean{⊔} is
the sum of subspaces (it is direct, which the dimension formula uses). \lean{hR} is ``no roots''.
For Leech, \lean{leechVOA.grading 2} is by definition \lean{gradingLatGen leech 2}, and the three
hypotheses are \lean{leech_even}, \lean{leech_posDef}, \lean{leech_no_roots}; the count $196560$ is
\lean{Leech.repCount_four}. The textbook statement and the Lean one agree. The result is
\emph{not} vacuous in the way a dimension count could be (e.g.\ by computing \lean{finrank} of an
infinite-dimensional space as $0$): the space is proved finite-dimensional and the number is
positive.
\end{dictionary}

\begin{physics}
The weight-two operators of the torus theory are of three types: the oscillator descendants $\partial^2X^a$, the bilinears $\partial X^a\partial X^b$ ($24+300$ of them), and the vertex operators $:\!e^{i\alpha\cdot X}\!:$ with $\tfrac12\langle\alpha,\alpha\rangle=2$, one for each of the $196560$ minimal vectors of $\Leech$. The total $196884=24+300+196560$ is the dimension of the weight-two space of $V_\Leech$, which numerically coincides with that of the Moonshine module's Griess algebra. This coincidence reflects that both characters equal $J+\text{const}$, and is not used in proofs.
\end{physics}

\paragraph{Proof.} A summand $(F_r)_{2-\frac12\langle\alpha,\alpha\rangle}\otimes e^\alpha$ is nonzero only if
$\tfrac12\langle\alpha,\alpha\rangle\le2$, i.e.\ $\langle\alpha,\alpha\rangle\in\{0,2,4\}$. Norm $0$ means
$\alpha=0$, giving $(F_r)_2\otimes e^0$; norm $2$ does not occur; norm $4$ gives
$(F_r)_0\otimes e^\alpha=\C e^\alpha$. Count dimensions with Proposition~\ref{ch04:prop:fockwt2}; for
$r=24$, $\binom{25}{2}=300$.\qed

\subsection{The \texorpdfstring{$\theta$}{theta}-split and the counting trick}

\begin{theorem}[Weight two of $V_L^+$]\label{ch04:thm:plustwo}
Let $L$ be even, positive-definite and rootless, and let $S$ be a set containing exactly one
vector from each pair $\{\alpha,-\alpha\}$ of norm-$4$ vectors. Then
\begin{gather*}
  (V_L)_2^{+}=\operatorname{span}\{h_a(-1)h_b(-1)e^0\}\oplus\bigoplus_{\alpha\in S}\C(e^\alpha+e^{-\alpha}),
  \\
  \dim(V_L)_2^+=\binom{r+1}{2}+\#S=\binom{r+1}{2}+\tfrac12\#\{\langle\alpha,\alpha\rangle=4\}.
\end{gather*}
For the Leech lattice, $\dim(V_\Leech^+)_2=300+98280=98580$.
\end{theorem}

\begin{leancode}
-- VOA/Moonshine/PlusWeightTwo.lean
noncomputable def plusTwo : Submodule ℂ (LatModuleGen L) :=
  gradingLatGen L 2 ⊓ LinearMap.ker ((thetaLat L).toLinearMap - LinearMap.id)
noncomputable def halfFour : Finset L.carrier :=
  (normFour L hPos).filter fun α => WellOrderingRel α (-α)
theorem card_normFour_eq : (normFour L hPos).card = 2 * (halfFour L hPos).card

theorem finrank_plusTwo (hEven : L.Even) (hR : ∀ v : L.carrier, L.form v v ≠ 2) :
    finrank ℂ (plusTwo L) = (rank L + 1).choose 2 + (halfFour L hPos).card

noncomputable def leechVOAPlusGradingTwoEquiv : leechVOAPlus.grading 2 ≃ₗ[ℂ] plusTwo leech
theorem leechVOAPlus_finrank_grading_two : finrank ℂ (leechVOAPlus.grading 2) = 98580
\end{leancode}

\begin{dictionary}
\lean{plusTwo L} is $(V_L)_2\cap\ker(\theta-1)$ (\lean{⊓} is intersection, \lean{LinearMap.ker} the
kernel), i.e.\ the $\theta$-fixed vectors of weight $2$ in $V_L$. \lean{halfFour L hPos} is the set
$S$: \lean{WellOrderingRel} is a well-order on $L$ supplied by the axiom of choice, and $S$ keeps
$\alpha$ when $\alpha$ precedes $-\alpha$; any choice of one vector per pair would do, and
\lean{card_normFour_eq} records $\#S=\tfrac12\#\{\langle\alpha,\alpha\rangle=4\}$. In
\lean{finrank_plusTwo}, \lean{hPos} is a section variable. The linear equivalence
\lean{leechVOAPlusGradingTwoEquiv} identifies the weight-$2$ space of the VOA \lean{leechVOAPlus}
(whose elements live in the subtype \lean{VPlus}) with the subspace \lean{plusTwo leech} of
$V_\Leech$; this is how a statement about the abstract VOA is reduced to the concrete subspace, and
it is what makes the final \lean{leechVOAPlus_finrank_grading_two} a statement about the actual
VOA $V_\Leech^+$ and not merely about a subspace defined by hand.
\end{dictionary}

\begin{physics}
The weight-two $\theta$-even states are the $300$ bilinears $\partial X^a\partial X^b$ (which include the stress tensor $\omega=\tfrac12\sum_a\partial X^a\partial X^a$) together with the $98280$ cosine combinations $:\!(e^{i\alpha\cdot X}+e^{-i\alpha\cdot X})\!:\;\sim\cos(\alpha\cdot X)$, one for each pair $\pm\alpha$ of minimal vectors. The $\theta$-odd ones, $\partial^2X^a$ and $\sin(\alpha\cdot X)$, are projected out. So $\dim(V_\Leech^+)_2=300+98280=98580$ is the untwisted contribution to the Griess algebra of $\Vnat$, whose remaining $98304=24\cdot2^{12}$ states come from the twisted sector. The counting trick is a pure linear algebra device with no direct physical meaning.
\end{physics}

\paragraph{Proof: the counting trick.} Write $D=\dim(V_L)_2=r+\binom{r+1}2+2\#S$
(Theorem~\ref{ch04:thm:wt2rootless} and $\#\{\langle\alpha,\alpha\rangle=4\}=2\#S$), and put
$V^\pm=(V_L)_2\cap\ker(\theta\mp1)$, the fixed and the odd subspace. The argument has four steps.

\emph{Step (i): $V^+\cap V^-=0$.} If $x\in V^+\cap V^-$ then $x=\theta x=-x$, so $2x=0$ and $x=0$.

\emph{Step (ii): an upper bound for the sum.} Since $V^+$ and $V^-$ meet only in $0$, their sum
$V^++V^-$ is direct, so $\dim V^++\dim V^-=\dim(V^+\oplus V^-)$. And $V^+\oplus V^-$ is a subspace of
$(V_L)_2$. Hence
\begin{equation}\label{ch04:eq:upper}
  \dim V^++\dim V^-\le D .
\end{equation}

\emph{Step (iii): a lower bound for each.} The quadratic states $h_a(-1)h_b(-1)e^0$ (there are
$\binom{r+1}2$, linearly independent) are fixed by $\theta$ (Proposition~\ref{ch04:prop:fockwt2}), and
for $\alpha\in S$ so is $e^\alpha+e^{-\alpha}=e^\alpha+\theta e^\alpha$. These $\binom{r+1}2+\#S$ vectors are
linearly independent: the $e^\alpha+e^{-\alpha}$ for distinct $\alpha\in S$ involve disjoint pairs of basis
vectors $e^{\pm\alpha}$ (that is why $S$ contains only one vector of each pair), and they live in
nonzero charge sectors, while the quadratic states live in charge $0$. So
$\dim V^+\ge\binom{r+1}2+\#S$. In the same way the $r$ states $h_a(-2)e^0$ and the $\#S$ vectors
$e^\alpha-e^{-\alpha}$ are independent and lie in $V^-$, so $\dim V^-\ge r+\#S$.

\emph{Step (iv): squeeze.} Adding the two lower bounds of (iii),
\begin{equation*}
  \dim V^++\dim V^-\ \ge\ \Big(\tbinom{r+1}2+\#S\Big)+\big(r+\#S\big)=D .
\end{equation*}
Together with \eqref{ch04:eq:upper}, $\dim V^++\dim V^-=D$. Now if either lower bound were strict,
say $\dim V^+>\binom{r+1}2+\#S$, then adding $\dim V^-\ge r+\#S$ would give
$\dim V^++\dim V^->D$, contradicting \eqref{ch04:eq:upper}. So both lower bounds are equalities:
$\dim V^+=\binom{r+1}2+\#S$ and $\dim V^-=r+\#S$. (As a by-product, the exhibited vectors are bases,
and $(V_L)_2=V^+\oplus V^-$.)\qed

The trick avoids computing a trace or a projection: we never need to show that the exhibited
vectors \emph{span} the fixed space, because the dimension count forces it.

\begin{example}[The weight-two spaces of $V_\Leech$, in a table]\label{ch04:ex:wt2}
\begin{center}
\begin{tabular}{lccc}
\toprule
states & number & $\theta$ & contributes to \\
\midrule
$h_a(-2)e^0$ & $24$ & $-1$ & $V^-$\\
$h_a(-1)h_b(-1)e^0$, $a\le b$ & $300$ & $+1$ & $V^+$\\
$e^\alpha+e^{-\alpha}$, $\alpha\in S$ & $98280$ & $+1$ & $V^+$\\
$e^\alpha-e^{-\alpha}$, $\alpha\in S$ & $98280$ & $-1$ & $V^-$\\
\midrule
total & $196884$ & & $98580+98304$\\
\bottomrule
\end{tabular}
\end{center}
Thus $\dim(V_\Leech^+)_2=98580$ and $\dim(V_\Leech^-)_2=98304$. The number $196884$ appears here as
$\dim(V_\Leech)_2$, \emph{before} any orbifolding; the Griess algebra $\Vnat_2$ has the same
dimension for a different reason: $98580+98304$ where now $98304=2^{12}\cdot24$ is the dimension of
the weight-$2$ part of the twisted sector, $(F^{T})_{1/2}\otimes T$ with $\dim T=2^{12}$ (Chapter~\ref{ch:08}).
That $\dim(V_\Leech^-)_2=24+98280$ coincides with $2^{12}\cdot24$ is a numerical observation
--- equivalently $\dim(V_\Leech)_2=\dim\Vnat_2$ --- and it is \emph{not} used anywhere in the
library; it reflects the fact that both $V_\Leech$ and $\Vnat$ have characters equal to $J$ plus a
constant (\S\ref{ch04:sec:gradeddim}).
\end{example}

\begin{status}
\lean{finrank_gradingLatGen_two}, \lean{finrank_plusTwo}: proved for every even positive-definite
rootless lattice. \lean{leechVOA_finrank_grading_two} ($196884$) and
\lean{leechVOAPlus_finrank_grading_two} ($98580$): proved, unconditional; \texttt{\#print axioms}
\texttt{[propext, Classical.choice, Quot.sound]}. The file \leanpath{LatticeGradedDim.lean} also
re-derives both numbers from the general formulas of the next section
(\lean{leechVOA_finrank_grading_two_theta}, \lean{leechVOAPlus_finrank_grading_two_theta}), an
internal consistency check.
\end{status}

%======================================================================
\section{The graded dimension of \texorpdfstring{$V_L$}{VL} and of \texorpdfstring{$V_L^+$}{VL+}}\label{ch04:sec:gradeddim}

\subsection{Coloured partitions}

\begin{definition}[Fock counts]\label{ch04:def:fockcount}
For $r,w\in\N$ let $P(r,w)$ be the number of monomials of weight $w$ in $F_r$ --- the number of
$r$-coloured partitions of $w$ --- and let $P^+(r,w)$, $P^-(r,w)$ be the numbers of those of even,
respectively odd, degree. The \emph{signed count} is $P^{\mathrm{s}}(r,w)=P^+(r,w)-P^-(r,w)$.
\end{definition}

\begin{leancode}
-- VOA/Examples/HeisenbergNVertex.lean
def weightDegSetN (r w : ℕ) : Set ((ℕ × Fin r) →₀ ℕ) :=
  {d : (ℕ × Fin r) →₀ ℕ | Finsupp.weight weightFnN d = w}

-- VOA/Moonshine/LatticeGradedDim.lean
def fockSetEven (w : ℕ) : Set ((ℕ × Fin r) →₀ ℕ) :=
  {d | Finsupp.weight weightFnN d = w ∧ Even d.degree}
def fockSetOdd (w : ℕ) : Set ((ℕ × Fin r) →₀ ℕ) :=
  {d | Finsupp.weight weightFnN d = w ∧ Odd d.degree}
noncomputable def fockCount (w : ℕ) : ℕ := (weightDegSetN r w).ncard
noncomputable def fockCountEven (w : ℕ) : ℕ := (fockSetEven r w).ncard
noncomputable def fockCountOdd (w : ℕ) : ℕ := (fockSetOdd r w).ncard
noncomputable def fockSigned (r w : ℕ) : ℤ := ∑ d ∈ fockFinset r w, (-1 : ℤ) ^ d.degree

theorem finrank_gradingWN (w : ℕ) : finrank ℂ (gradingWN r w) = fockCount r w
theorem fockSigned_eq (r w : ℕ) :
    fockSigned r w = (fockCountEven r w : ℤ) - fockCountOdd r w
\end{leancode}

\begin{dictionary}
\lean{Set.ncard} is the cardinality of a set as a natural number (and $0$ for infinite sets ---
again a potential source of vacuity; here the sets are finite, \lean{weightDegSetN_finite}).
\lean{fockCount r w}, \lean{fockCountEven}, \lean{fockCountOdd} are $P$, $P^+$, $P^-$;
\lean{fockFinset r w} is the same set of exponents as a \lean{Finset}, and
\lean{fockSigned r w} $=\sum_d(-1)^{\deg d}=P^{\mathrm s}(r,w)$. \lean{finrank_gradingWN} says that
the monomials of weight $w$ form a basis of $(F_r)_w$. Since $\sigma(x^d)=(-1)^{\deg d}x^d$,
$P^{\mathrm s}(r,w)$ is the trace of $\sigma$ on $(F_r)_w$; the library never forms this trace and
works with the counts directly.
\end{dictionary}

\begin{physics}
$P(r,w)$ is the number of states at energy $w$ of $r$ free chiral bosons, and $P^+,P^-$ separate them by $\mathbb Z_2$ parity (even or odd number of quanta). The low values $1,\,r,\,r+\binom{r+1}2,\dots$ are the first terms of the partition function $q^{r/24}\eta(q)^{-r}$ of the bosons. For $r=24$ they give $1,24,324,3200,\dots$, the coefficients of $\eta^{-24}$, the familiar spectrum of the bosonic string.
\end{physics}

\begin{proposition}[Small weights]\label{ch04:prop:smallwt}
$P(r,0)=1$, $P(r,1)=r$, $P(r,2)=r+\binom{r+1}2$, $P(r,3)=r+r^2+\binom{r+2}3$; and
$P^-(r,2)=r$, $P^+(r,2)=\binom{r+1}2$, $P^+(r,3)=r^2$, $P^-(r,3)=r+\binom{r+2}3$.
\end{proposition}

\begin{leancode}
-- VOA/Moonshine/LatticeGradedDim.lean
theorem fockCount_zero : fockCount r 0 = 1
theorem fockCount_one : fockCount r 1 = r
theorem fockCount_two : fockCount r 2 = r + (r + 1).choose 2
theorem fockCountOdd_two : fockCountOdd r 2 = r
theorem fockCountEven_two : fockCountEven r 2 = (r + 1).choose 2
theorem fockCountEven_three : fockCountEven r 3 = r * r
theorem fockCountOdd_three : fockCountOdd r 3 = r + (r + 2).choose 3
theorem fockCount_three : fockCount r 3 = r + r * r + (r + 2).choose 3
\end{leancode}

\begin{physics}
These are the oscillator excitations at level $3$ of $r$ bosons: a quantum of energy $3$, a quantum of energy $2$ together with one of energy $1$, or three quanta of energy $1$. The parity counts $P^\pm$ record how the $\mathbb Z_2$ symmetry acts; the difference $P^+-P^-=\operatorname{tr}\theta$ is the term that will appear in the twisted partition function.
\end{physics}

\paragraph{Proof.} Enumerate the monomials. Weight $3$: $h_a(-3)$ ($r$, degree $1$, odd),
$h_a(-2)h_b(-1)$ ($r^2$, ordered pairs since the two oscillators have different weights; degree $2$,
even), $h_a(-1)h_b(-1)h_c(-1)$ (multisets of size $3$: $\binom{r+2}3$; degree $3$, odd).\qed

For $r=24$: $P(24,3)=24+576+2600=3200$, $P^+(24,3)=576$, $P^-(24,3)=2624$.

\subsection{Generating functions}

\begin{theorem}[Fock generating functions]\label{ch04:thm:genfun}
In the ring $\C[[q]]$ of formal power series,
\begin{equation*}
  \sum_{w\ge0}P(r,w)q^w=\prod_{n\ge1}(1-q^n)^{-r},
  \qquad
  \sum_{w\ge0}\big(P^+(r,w)-P^-(r,w)\big)q^w=\prod_{n\ge1}(1+q^n)^{-r},
\end{equation*}
the infinite products converging in the $q$-adic (coefficientwise) topology.
\end{theorem}

\begin{leancode}
-- VOA/Moonshine/LatticeGradedDim.lean
noncomputable def fockGenFun (r : ℕ) : PowerSeries ℂ := PowerSeries.mk fun w => (fockCount r w : ℂ)
noncomputable def fockSignedGenFun (r : ℕ) : PowerSeries ℂ :=
  PowerSeries.mk fun w => (fockSigned r w : ℂ)

open scoped PowerSeries.WithPiTopology in
theorem hasProd_fockGenFun (r : ℕ) :
    HasProd (fun i : ℕ => ((1 - (PowerSeries.X : ℂ⟦X⟧) ^ (i + 1))⁻¹) ^ r) (fockGenFun r)

open scoped PowerSeries.WithPiTopology in
theorem hasProd_fockSignedGenFun (r : ℕ) :
    HasProd (fun i : ℕ => ((1 + (PowerSeries.X : ℂ⟦X⟧) ^ (i + 1))⁻¹) ^ r)
      (fockSignedGenFun r)
\end{leancode}

\begin{dictionary}
\lean{PowerSeries ℂ} $=$ \lean{ℂ⟦X⟧} is $\C[[q]]$ with the variable called $X$;
\lean{PowerSeries.mk f} is $\sum_w f(w)X^w$. \lean{⁻¹} is the inverse in $\C[[q]]$ (which exists for
$1\mp X^{i+1}$, constant term $1$). \lean{HasProd f a} is Mathlib's statement that the infinite
product $\prod_{i\in\N}f(i)$ converges (unconditionally, along finite subsets) to $a$, for the
topology fixed by \lean{open scoped PowerSeries.WithPiTopology}: the product topology on the
coefficient sequence, i.e.\ coefficientwise convergence where each coefficient is eventually
constant. This is the standard meaning of an infinite product of formal power series. The index
shift $i\mapsto i+1$ makes $n=i+1$ run over $n\ge1$. The statements are exactly
Theorem~\ref{ch04:thm:genfun}.
\end{dictionary}

\begin{physics}
The first identity is the partition function of $r$ free chiral bosons, $Z_{\mathrm{osc}}=\operatorname{tr}\,q^{L_0}=\prod_n(1-q^n)^{-r}$ (each mode $\alpha_{-n}$ is a harmonic oscillator of energy $n$, contributing $(1-q^n)^{-1}$). The second is the same trace with the reflection $\theta$ inserted, $\operatorname{tr}\,\theta\,q^{L_0}=\prod_n(1+q^n)^{-r}$, since each excitation then contributes $\sum_j(-q^n)^j$. These traces are the building blocks of the orbifold partition function $\tfrac12(Z+\operatorname{tr}\theta q^{L_0})$.
\end{physics}

\paragraph{Proof.} For one colour, $\prod_{n\ge1}(1-q^n)^{-1}=\prod_n\sum_{j\ge0}q^{nj}$ is Euler's
generating function of partitions: choosing $j_n$ parts equal to $n$ for each $n$. A one-colour
monomial $\prod_m x_{(m,0)}^{d_m}$ is exactly the partition with $d_m$ parts equal to $m+1$, so
$\sum_wP(1,w)q^w$ is Euler's series. For $r$ colours, a monomial is an $r$-tuple of one-colour
monomials and weights add, so the generating function is the $r$-th power. For the signed count,
attach $(-1)^{j}$ to $j$ parts equal to $n$: $\sum_j(-1)^jq^{nj}=(1+q^n)^{-1}$, and the total sign is
$(-1)^{\deg}$ because the degree is the number of parts. In Lean the one-colour case is reduced to
Mathlib's partition generating function by an explicit bijection between monomials and
partitions.\qed

\begin{example}[Low coefficients for $r=24$]\label{ch04:ex:fock24}
\begin{equation*}
  \prod_{n\ge1}(1-q^n)^{-24}=1+24q+324q^2+3200q^3+\cdots,\qquad
  \prod_{n\ge1}(1+q^n)^{-24}=1-24q+276q^2-2048q^3+\cdots.
\end{equation*}
Check against Proposition~\ref{ch04:prop:smallwt}: $324=24+300$, $276=300-24$, $3200=24+576+2600$,
$-2048=576-2624$. (Directly: the $q^3$ coefficient of $\prod(1+q^n)^{-24}$ is
$-\binom{26}3+24\cdot24-24=-2600+576-24=-2048$.)
\end{example}

\subsection{The dimension formulas}

\begin{theorem}[Graded dimension of $V_L$]\label{ch04:thm:dimVL}
Let $L$ be even and positive-definite of rank $r$. For every $k\in\N$,
\begin{equation*}
  \dim(V_L)_k=\sum_{n=0}^{k}\#\{\alpha\in L:\langle\alpha,\alpha\rangle=2n\}\cdot P(r,k-n),
\end{equation*}
the coefficient of $q^k$ in $\Theta_L(q)\prod_{n\ge1}(1-q^n)^{-r}$, where
$\Theta_L(q)=\sum_{\alpha\in L}q^{\langle\alpha,\alpha\rangle/2}$ is the theta series of $L$.
\end{theorem}

\begin{leancode}
-- VOA/Moonshine/LatticeGradedDim.lean
noncomputable def normEq (m : ℤ) : Finset L.carrier := (normEq_finite L hPos m).toFinset
theorem mem_normEq {m : ℤ} {α : L.carrier} : α ∈ normEq L hPos m ↔ L.form α α = m

theorem finrank_gradingLatGen_eq_sum_theta (hEven : L.Even) (k : ℕ) :
    finrank ℂ (gradingLatGen L k)
      = ∑ n ∈ Finset.range (k + 1), (normEq L hPos (2 * n)).card * fockCount (rank L) (k - n)
\end{leancode}

\begin{dictionary}
Section variables \lean{(L : IntegralLattice) (hPos : L.PosDef)}. \lean{normEq L hPos m} is the
finite set of vectors of norm $m$. \lean{Finset.range (k + 1)} is $\{0,\dots,k\}$; \lean{k - n} is
natural-number subtraction, which is honest here since $n\le k$. \lean{gradingLatGen L k} has
\lean{k : ℕ} coerced to $\Z$. The formula is Theorem~\ref{ch04:thm:dimVL} verbatim; the identification
with the $q^k$ coefficient of the product series is the Cauchy product of
Theorem~\ref{ch04:thm:genfun} with $\Theta_L$, made in Chapter~\ref{ch:05} where characters are
treated as power series.
\end{dictionary}

\begin{physics}
This is the partition function of the compact boson as a lattice sum: $\operatorname{tr}\,q^{L_0-r/24}=\Theta_L(q)/\eta(q)^r$, with the momentum sum $\Theta_L$ over the lattice times the oscillator contribution. For the Leech lattice, $\Theta_\Leech/\eta^{24}=J+24$ up to the usual $q^{-1}$ shift, i.e.\ $q^{-1}+24+196884\,q+\cdots$; the constant $24$ is the $24$ currents, to be removed by the orbifold.
\end{physics}

\paragraph{Proof.} $(V_L)_k=\bigoplus_\alpha(F_r)_{k-\frac12\langle\alpha,\alpha\rangle}\otimes e^\alpha$, the
sum being over the finitely many $\alpha$ with $\tfrac12\langle\alpha,\alpha\rangle\le k$; the summand has
dimension $P(r,k-\tfrac12\langle\alpha,\alpha\rangle)$. Group the $\alpha$ by $n=\tfrac12\langle\alpha,\alpha\rangle$.\qed

\begin{theorem}[Graded dimension of $V_L^+$]\label{ch04:thm:dimVLplus}
With $L$ as above, let $(V_L)_k^\pm$ be the $\pm1$-eigenspaces of $\theta$ on $(V_L)_k$. Then
\begin{equation*}
  2\dim(V_L)_k^+=\dim(V_L)_k+P^+(r,k)-P^-(r,k),
\end{equation*}
i.e.\ $\dim(V_L)_k^+=\tfrac12\big(\dim(V_L)_k+\tr(\theta|(V_L)_k)\big)$ with
$\tr(\theta|(V_L)_k)=P^{\mathrm s}(r,k)$, the coefficient of $q^k$ in $\prod_{n\ge1}(1+q^n)^{-r}$.
Explicitly, with $S$ a set of representatives of the pairs $\pm\alpha\neq0$ of half-norm $\le k$,
\begin{equation*}
  \dim(V_L)_k^\pm=P^\pm(r,k)+\sum_{\alpha\in S}P\big(r,k-\tfrac12\langle\alpha,\alpha\rangle\big).
\end{equation*}
\end{theorem}

\begin{leancode}
-- VOA/Moonshine/LatticeGradedDim.lean
noncomputable def plusK (k : ℤ) : Submodule ℂ (LatModuleGen L) :=
  gradingLatGen L k ⊓ LinearMap.ker ((thetaLat L).toLinearMap - LinearMap.id)
noncomputable def minusK (k : ℤ) : Submodule ℂ (LatModuleGen L) :=
  gradingLatGen L k ⊓ LinearMap.ker ((thetaLat L).toLinearMap + LinearMap.id)
noncomputable def halfLE (k : ℤ) : Finset L.carrier :=
  (normLE L hPos k).filter fun α => WellOrderingRel α (-α)

theorem finrank_plusK_minusK (hEven : L.Even) (k : ℕ) :
    finrank ℂ (plusK L k) = fockCountEven (rank L) k
        + ∑ α ∈ halfLE L hPos k, fockCount (rank L) ((k : ℤ) - halfNorm L α).toNat ∧
      finrank ℂ (minusK L k) = fockCountOdd (rank L) k
        + ∑ α ∈ halfLE L hPos k, fockCount (rank L) ((k : ℤ) - halfNorm L α).toNat

theorem two_mul_finrank_plusK_eq_signed (L : IntegralLattice) (hPos : L.PosDef) (hEven : L.Even)
    (k : ℕ) :
    2 * (finrank ℂ (plusK L k) : ℤ) = finrank ℂ (gradingLatGen L k) + fockSigned (rank L) k

noncomputable def leechVOAPlusGradingEquiv (k : ℤ) : leechVOAPlus.grading k ≃ₗ[ℂ] plusK leech k
\end{leancode}

\begin{dictionary}
\lean{plusK L k}, \lean{minusK L k} are $(V_L)^\pm_k$. \lean{normLE L hPos k} is the finite set
$\{\alpha:\tfrac12\langle\alpha,\alpha\rangle\le k\}$ and \lean{halfLE} keeps one of each pair $\pm\alpha$
(it excludes $\alpha=0$, which is its own negative). \lean{.toNat} truncates an integer to $\N$;
harmless because $k-\tfrac12\langle\alpha,\alpha\rangle\ge0$ for $\alpha\in$ \lean{normLE}. The equation
\lean{two_mul_finrank_plusK_eq_signed} is stated in $\Z$ (the cast \lean{(… : ℤ)}) because
$P^{\mathrm s}$ can be negative. \lean{leechVOAPlusGradingEquiv k} identifies the weight-$k$ space
of the VOA $V_\Leech^+$ with \lean{plusK leech k}, exactly as in weight $2$.
\end{dictionary}

\begin{physics}
The dimension of the invariant subspace at each level is $\tfrac12(\dim+\operatorname{tr}\theta)$, i.e.\ the projection onto the untwisted $\mathbb Z_2$-even sector $\tfrac12(Z+\operatorname{tr}\theta q^{L_0})$ of an orbifold partition function. Only the neutral (zero-momentum) sector contributes to $\operatorname{tr}\theta$, because $\theta$ exchanges every charged sector $\alpha\ne0$ with its partner $-\alpha$, so charged states are off-diagonal and have zero trace. In this way the formula keeps the $98280$ pairs of cosine states at weight $2$ and kills the sine ones.
\end{physics}

\paragraph{Proof.} The counting trick of Theorem~\ref{ch04:thm:plustwo}, in every weight. In charge
$0$, $\theta=\sigma\otimes1$ acts on monomials by $(-1)^{\deg}$, so $(F_r)_k^{\mathrm{even}}\otimes e^0$
is fixed and $(F_r)_k^{\mathrm{odd}}\otimes e^0$ is anti-fixed. For $\alpha\in S$, the map
$x\mapsto x\pm\theta x$ is injective on $(F_r)_{k-\frac12\langle\alpha,\alpha\rangle}\otimes e^\alpha$ (its image
has a component in the sector $\alpha$, where $\theta x$ has none) and lands in $(V_L)^\pm_k$. These
lower bounds for $\dim(V_L)^+_k$ and $\dim(V_L)^-_k$ add up to $P(r,k)+2\sum_{\alpha\in S}P(\dots)=\dim(V_L)_k$,
so both are equalities. Subtracting the two equalities gives the trace form. Physically: the
nonzero charge sectors come in pairs $\{\alpha,-\alpha\}$ swapped by $\theta$ and contribute nothing to
the trace; only the neutral sector, where $\theta=(-1)^{\#\text{oscillators}}$, contributes.\qed

\begin{theorem}[Weight three]\label{ch04:thm:wt3}
$\dim(V_\Leech)_3=21493760$ and $\dim(V_\Leech^+)_3=10745856$.
\end{theorem}

\begin{leancode}
-- VOA/Moonshine/LatticeGradedDim.lean
theorem leechVOA_finrank_grading_three : finrank ℂ (leechVOA.grading 3) = 21493760
theorem leechVOAPlus_finrank_grading_three : finrank ℂ (leechVOAPlus.grading 3) = 10745856
\end{leancode}

\begin{dictionary}
No hypotheses; the inputs are the Leech counts $\#\{\langle\alpha,\alpha\rangle=0,2,4,6\}=1,0,196560,16773120$
(\lean{leech_normEq_zero_card}, \dots, \lean{leech_normEq_six_card}, from
\lean{Leech.repCount_four} and \lean{Leech.repCount_six}).
\end{dictionary}

\begin{physics}
This is the next level of the same untwisted orbifold sector: $10745856$ states at energy $3$ survive the $\mathbb Z_2$ projection out of $21493760$. Adding the twisted-sector count of Chapter~\ref{ch:08} gives the full weight-three Moonshine space and matches the coefficient $21493760$ of $q^2$ in $J(q)$. Heuristic, as a consistency check on the physics of the orbifold.
\end{physics}

\paragraph{Proof.} By Theorem~\ref{ch04:thm:dimVL} with $r=24$:
\begin{align*}
  \dim(V_\Leech)_3&=1\cdot P(24,3)+0\cdot P(24,2)+196560\cdot P(24,1)+16773120\cdot P(24,0)\\
  &=3200+4717440+16773120=21493760 .
\end{align*}
By Theorem~\ref{ch04:thm:dimVLplus}, $2\dim(V_\Leech^+)_3=21493760+576-2624=21491712$, so
$\dim(V_\Leech^+)_3=10745856$.\qed

\begin{example}[The character of $V_\Leech$ and its twin]\label{ch04:ex:char}
Theorem~\ref{ch04:thm:dimVL} gives
\begin{equation*}
  \sum_k\dim(V_\Leech)_kq^{k-1}=\frac{\Theta_\Leech(q)}{q\prod_n(1-q^n)^{24}}
  =q^{-1}+24+196884q+21493760q^2+\cdots=J(q)+24,
\end{equation*}
using $\Theta_\Leech=1+196560q^2+16773120q^3+\cdots$ (in the variable $q^{\langle\alpha,\alpha\rangle/2}$,
the norm-$4$ vectors at $q^2$). For the $\theta$-trace, $\sum_k\tr(\theta|(V_\Leech)_k)q^{k-1}
=q^{-1}\prod_n(1+q^n)^{-24}=q^{-1}-24+276q-2048q^2+\cdots$. Hence
\begin{equation*}
  \sum_k\dim(V_\Leech^+)_kq^{k-1}=q^{-1}+0+98580q+10745856q^2+\cdots .
\end{equation*}
Adding the twisted contribution $98304q+10747904q^2+\cdots$ of Chapter~\ref{ch:08} gives
$q^{-1}+196884q+21493760q^2+\cdots=J(q)$: the miracle that the constant term $24$ is removed and
nothing else changes. The full identity for every coefficient is proved in Chapter~\ref{ch:05}.
\end{example}

\begin{status}
\lean{finrank_gradingLatGen_eq_sum_theta}, \lean{finrank_plusK_minusK},
\lean{two_mul_finrank_plusK_eq_signed}, \lean{hasProd_fockGenFun},
\lean{hasProd_fockSignedGenFun}: proved for every even positive-definite $L$ (and every $r$).
\lean{leechVOA_finrank_grading_three}, \lean{leechVOAPlus_finrank_grading_three}: proved,
unconditional. \texttt{\#print axioms} for all of these: \texttt{[propext, Classical.choice,
Quot.sound]}. (The weight-$3$ proofs use the kernel-checked \lean{decide} to evaluate
$\binom{26}{3}=2600$, not \lean{native_decide}.)
\end{status}
%======================================================================
\section{Generation: \texorpdfstring{$V_\Lambda^+$}{V+} is generated by the \texorpdfstring{$e^\beta+e^{-\beta}$}{e(b)+e(-b)}}\label{ch04:sec:generation}

\subsection{Why generation is needed}

Several later constructions define a structure (an invariant bilinear form, an intertwining
operator) and must check a compatibility with the action of \emph{every} $u\in V_\Leech^+$. Checking
it directly for every $u$ is hopeless; checking it on a generating set is easy, provided one knows
that the set of $u$ for which the compatibility holds is closed under $n$-th products. The form
of \S\ref{ch04:sec:form} below and the form \lean{TPForm} on the twisted sector
(Chapter~\ref{ch:09}) are both handled this way. What is needed is therefore a precise generation
theorem.

\begin{definition}\label{ch04:def:xbeta}
For $\beta\in\Leech$ put $x_\beta=e^\beta+e^{-\beta}\in V_\Leech^+$. (It is $\theta$-fixed since
$\theta e^\beta=e^{-\beta}$; note $x_0=2e^0$ and $x_{-\beta}=x_\beta$.)
\end{definition}

\begin{leancode}
-- VOA/Moonshine/CyclicFrameSectors.lean
noncomputable def frameX0 (β : leech.carrier) : LatModuleGen leech :=
  latChargeGen leech β + latChargeGen leech (-β)
noncomputable def frameX0P (β : leech.carrier) : VPlus := ⟨frameX0 β, frameX0_mem_plus β⟩
\end{leancode}

\begin{dictionary}
\lean{frameX0 β} is $x_\beta$ as a vector of $V_\Leech$; \lean{frameX0P β} is the same vector as an
element of $V_\Leech^+$, i.e.\ paired (\lean{⟨…, …⟩}) with the proof \lean{frameX0_mem_plus β} that it
is $\theta$-fixed. (The name comes from the Golay-frame chapter, where $x_\beta$ with special $\beta$
is used as a cyclic vector.)
\end{dictionary}

\begin{physics}
The operator $x_\beta=e^\beta+e^{-\beta}\sim2\cos(\beta\cdot X)$ is the $\mathbb Z_2$-even combination of a vertex operator and its reflection: the simplest local operators that survive the orbifold projection with nonzero momentum. For $\beta$ of norm $4$ these are the weight-two cosine operators of the Griess algebra; for general $\beta$ they have higher weight $\tfrac12\langle\beta,\beta\rangle$.
\end{physics}

\begin{theorem}[Generation]\label{ch04:thm:gen}
Let $S\subseteq V_\Leech^+$ be a linear subspace such that $\chfourone\in S$, $u_{(n)}v\in S$ for all
$u,v\in S$ and $n\in\Z$, and $x_\beta\in S$ for every $\beta\in\Leech$. Then $S=V_\Leech^+$.
\end{theorem}

\begin{leancode}
-- VOA/Moonshine/Generation.lean
theorem leechPlus_generated : ∀ S : Submodule ℂ VPlus, leechVOAPlus.vacuum ∈ S →
    (∀ u ∈ S, ∀ v ∈ S, ∀ n : ℤ, leechVOAPlus.smul u n v ∈ S) →
    (∀ β : leech.carrier, frameX0P β ∈ S) → S = ⊤
\end{leancode}

\begin{dictionary}
\lean{S : Submodule ℂ VPlus} is a linear subspace of $V_\Leech^+$. The three hypotheses are, in
order: $\chfourone\in S$; closure under all $n$-th products of $V_\Leech^+$ (\lean{leechVOAPlus.smul u n v}
$=u_{(n)}v$, computed in $V_\Leech$ since the sub-VOA has the restricted vertex map);
$x_\beta\in S$ for all $\beta$. The conclusion \lean{S = ⊤} is $S=V_\Leech^+$ (\lean{⊤} is the whole
space). This is the standard meaning of ``the $x_\beta$ generate $V_\Leech^+$ as a vertex algebra''.
Note: generation is in the sense of \emph{$n$-th products}, for \emph{all} $\beta\in\Leech$ (not only
the minimal ones), and the vacuum hypothesis is redundant ($x_0=2\chfourone$). The statement could
have been weak in two ways, neither of which occurs: it is not restricted to a subset of
$V_\Leech^+$, and it does not assume $S$ is a VOA (only a subspace).
\end{dictionary}

\begin{physics}
The theorem says that the chiral algebra of the orbifold is the OPE closure of the cosine operators: every $\theta$-even operator, including the bilinears $\partial X^a\partial X^b$ and the stress tensor, arises by repeatedly taking OPE coefficients (modes $u_{(n)}v$) of the cosines $\cos(\beta\cdot X)$. Heuristically: the OPE of $\cos(\beta\cdot X)$ with itself produces the $\partial X\partial X$ terms in its regular part. This closure property is what allows checking invariance (of a bilinear form, of an intertwiner) on a small set of generating fields rather than on all of $V_\Leech^+$.
\end{physics}

\begin{warning}
Generation by $n$-th products is \emph{weaker} than what is needed in another part of the project.
The Zhu-algebra arguments of Chapter~\ref{ch:12} need \emph{Zhu generation}. Recall (it is
defined and discussed there) that Zhu's algebra is $A(V)=V/O(V)$ with an associative product $a*b$
built from the modes of $Y(a,z)$; a set $G$ is \emph{Zhu-generating} if every subspace containing
$\chfourone$, $O(V)$ and $G$ and closed under $*$ is all of $V$ (library: \lean{ZhuGenerated} in
\leanpath{VOA/Moonshine/SigmaExistZhu.lean}). This closes under a single specific bilinear product
modulo $O(V)$, not under all $n$-th products, so neither notion implies the other in an obvious way.
The project recorded as a dead end that weak ($n$-th product) generation does NOT
reduce the top-level clause there to generators. Theorem~\ref{ch04:thm:gen} is used only for the
invariance arguments of \S\ref{ch04:sec:form} and Chapter~\ref{ch:09}, where $n$-th product
generation is exactly the right notion.
\end{warning}

\subsection{The proof}

The proof is elementary and works for every even positive-definite lattice; only the names are
Leech-specific in the library. Let $P=\tfrac12(1+\theta)$ be the projection of
Remark~\ref{ch04:rem:proj} and let $G=P^{-1}(S)\subseteq V_\Leech$. Since $P$ is the identity on
$V_\Leech^+$, it suffices to show $G=V_\Leech$, i.e.\ that every $p\otimes e^\gamma$ lies in $G$. We
show this by induction on the Fock weight $w$ of $p$, for all charges $\gamma$ at once.

\emph{Base.} $w=0$: $P(e^\gamma)=\tfrac12x_\gamma\in S$.

\emph{Step 1: the safe step --- why it is needed.} The only tool we have for producing new
elements of $S$ is the product: if $s\in S$ then $x_{a\,(n)}s\in S$. We would like to use this to
show that $e^a_{(n)}v\in G$ whenever $v\in G$, because the operators $e^a_{(n)}$ on $V_\Leech$ are
completely explicit, by \eqref{ch04:eq:Yealpha}. The difficulty is that $x_a=e^a+e^{-a}$ is a
\emph{sum}, and $P(v)$ is a sum too, so $x_{a\,(n)}P(v)$ has four terms of which we only want one
(and its $\theta$-mirror image). The safe step is the observation that for $n$ large enough the
two unwanted ``cross terms'' vanish identically, so that $x_{a\,(n)}P(v)$ is exactly $P(e^a_{(n)}v)$.

\emph{When does $e^b_{(n)}$ kill a vector?} From \eqref{ch04:eq:Yealpha}, for $q$ of Fock weight $w$
and any $b,c\in\Leech$,
\begin{equation}\label{ch04:eq:eamode}
  e^b_{(n)}(q\otimes e^c)=\varepsilon(b,c)\,EE(b)_{t}\,q\otimes e^{b+c},\qquad t=-n-1-\langle b,c\rangle,
\end{equation}
where $EE(b)_t$ is the coefficient of $z^t$ in $EE(b,z)=E^-(b,z)E^+(b,z)$. (The mode $u_{(n)}$ is the
coefficient of $z^{-n-1}$, and the factor $z^{\langle b,c\rangle}$ shifts the index.) Multiplying out,
$EE(b)_t=\sum_{N-M=t,\ N,M\ge0}S_N(b)\,E^+(b)_{-M}$. The operator $E^+(b)_{-M}$ is a polynomial in
annihilation operators of total weight $M$, so it lowers the Fock weight by exactly $M$ and kills
$q$ if $M>w$. Thus only terms with $M\le w$ survive, and for them $t=N-M\ge-M\ge-w$. In other words,
\begin{equation}\label{ch04:eq:kill}
  t<-w\ \Longrightarrow\ EE(b)_tq=0,\qquad\text{i.e.}\qquad
  n\ge w-\langle b,c\rangle\ \Longrightarrow\ e^b_{(n)}(q\otimes e^c)=0
\end{equation}
(the two conditions agree: $-n-1-\langle b,c\rangle<-w\iff n>w-1-\langle b,c\rangle\iff n\ge w-\langle b,c\rangle$).
Physically this is the statement that the operator product $e^b(z)\,(q\otimes e^c)(0)$ has a pole of
order at most $w-\langle b,c\rangle$: the momenta contribute $z^{\langle b,c\rangle}$ and the oscillators of $q$
can lower this by at most $w$.

\begin{lemma}[Safe step]\label{ch04:lem:safe}
If $q\in(F_{24})_w$ and $n\ge w+\langle a,c\rangle$, then
$\;x_{a\,(n)}\,P(q\otimes e^c)=P\big(e^a_{(n)}(q\otimes e^c)\big)$.
\end{lemma}

\begin{leancode}
-- VOA/Moonshine/Generation.lean
theorem plusProj_smul (a c : leech.carrier) {n : ℤ} {w : ℕ} {q : FockN (rank leech)}
    (hq : q ∈ gradingWN (rank leech) w) (hn : (w : ℤ) + leech.form a c ≤ n) :
    plusProj (leechVOA.smul (latChargeGen leech a) n (latStateGen leech q c))
      = leechVOAPlus.smul (frameX0P a) n (plusProj (latStateGen leech q c))
\end{leancode}

\begin{dictionary}
\lean{plusProj} is $P$; the left side is $P(e^a_{(n)}(q\otimes e^c))$, computed in $V_\Leech$; the
right side is $x_{a\,(n)}P(q\otimes e^c)$, computed in $V_\Leech^+$. The hypotheses are $q\in(F_{24})_w$
and $n\ge w+\langle a,c\rangle$.
\end{dictionary}

\begin{physics}
This lemma is a statement about when the OPE of a cosine $x_a=e^a+e^{-a}$ with a projected state keeps only the ``direct'' term: for a high enough mode number $n$ the cross terms, in which the momenta add rather than cancel, have an OPE with too weak a pole to appear, since $e^b(z)\,(q\otimes e^c)(0)$ has a pole of order at most $w-\langle b,c\rangle$. It is a pole-order estimate for the OPE of lattice vertex operators and helps one extract a single exponential from a cosine. Technical lemma.
\end{physics}

\paragraph{Proof.} Expand $x_{a\,(n)}\cdot2P(q\otimes e^c)=(e^a+e^{-a})_{(n)}(q\otimes e^c+\sigma q\otimes
e^{-c})$ into four terms. The cross term $e^{-a}_{(n)}(q\otimes e^c)$ is \eqref{ch04:eq:kill} with
$b=-a$: it vanishes when $n\ge w-\langle -a,c\rangle=w+\langle a,c\rangle$. The cross term
$e^{a}_{(n)}(\sigma q\otimes e^{-c})$ is \eqref{ch04:eq:kill} with $b=a$, charge $-c$ and the
weight-$w$ vector $\sigma q$: it vanishes when $n\ge w-\langle a,-c\rangle=w+\langle a,c\rangle$. Both
conditions are the hypothesis of the lemma. The remaining two are
$e^a_{(n)}(q\otimes e^c)$ and its image under $\theta$ (Theorem~\ref{ch04:thm:thetamode}:
$\theta(e^a_{(n)}(q\otimes e^c))=e^{-a}_{(n)}(\sigma q\otimes e^{-c})$). Their sum is
$2P(e^a_{(n)}(q\otimes e^c))$.\qed

Consequently: \emph{if $q\otimes e^c\in G$ and $n\ge w+\langle a,c\rangle$, then
$e^a_{(n)}(q\otimes e^c)\in G$} (because $P$ of it is $x_{a(n)}$ applied to an element of $S$).

\begin{example}[The safe step in rank one]\label{ch04:ex:saferank1}
Let $L=\Z\alpha$ with $\langle\alpha,\alpha\rangle=4$ (rank one, rootless), and take $a=\alpha$, $c=-\alpha$,
$q=1$ (so $w=0$). Then $P(e^{-\alpha})=\tfrac12x_\alpha$, and the lemma's condition is
$n\ge0+\langle\alpha,-\alpha\rangle=-4$. Compute the four terms of $x_{\alpha\,(n)}x_\alpha$ with
\eqref{ch04:eq:eamode}, where $t=-n-1-\langle b,c\rangle$:
\begin{center}
\begin{tabular}{llll}
\toprule
term & $\langle b,c\rangle$ & $t$ & value\\
\midrule
$e^{\alpha}_{(n)}e^{-\alpha}$ (wanted) & $-4$ & $3-n$ & $\varepsilon\,S_{3-n}(\alpha)\otimes e^0$\\
$e^{-\alpha}_{(n)}e^{\alpha}$ (its $\theta$-image) & $-4$ & $3-n$ & $\varepsilon\,S_{3-n}(-\alpha)\otimes e^0$\\
$e^{\alpha}_{(n)}e^{\alpha}$ (cross) & $4$ & $-n-5$ & $\varepsilon\,S_{-n-5}(\alpha)\otimes e^{2\alpha}$\\
$e^{-\alpha}_{(n)}e^{-\alpha}$ (cross) & $4$ & $-n-5$ & $\varepsilon\,S_{-n-5}(-\alpha)\otimes e^{-2\alpha}$\\
\bottomrule
\end{tabular}
\end{center}
(Here $q=1$, so $E^+$ acts trivially and $EE(b)_t1=S_t(b)$, which is $0$ for $t<0$; $\varepsilon$ denotes
the relevant sign $\pm1$.) For $n\ge-4$ the cross terms have $t=-n-5<0$ and vanish, so
$x_{\alpha\,(n)}x_\alpha=\varepsilon\,(S_{3-n}(\alpha)+S_{3-n}(-\alpha))\otimes e^0=2\varepsilon P(S_{3-n}(\alpha)\otimes e^0)$:
exactly the lemma. For $n=1$ this is $\varepsilon\,\alpha(-1)^2\chfourone$, as in Example~\ref{ch04:ex:genwt2}.
For the unsafe value $n=-5$, however, the cross terms are $\varepsilon\,(e^{2\alpha}+e^{-2\alpha})$: the product
is still an element of $S$, but it is \emph{no longer} $P$ of $e^\alpha_{(-5)}e^{-\alpha}$, and we learn
nothing about the latter. This is the whole content of the lemma, and every entry of the table can
be checked from \eqref{ch04:eq:eamode} and $S_N(-\alpha)=\sigma S_N(\alpha)$ alone.
\end{example}

\emph{Step 2: leading terms.} Fix a target weight $w\ge1$ and assume (induction hypothesis) that
all $p\otimes e^c$ with Fock weight $<w$ lie in $G$, for all $c$. For $d\ge0$ and $q$ of weight $k$,
\begin{equation}\label{ch04:eq:EEexpand}
  EE(a)_d\,q=\sum_{s=0}^{k}S_{d+s}(a)\,E^+(a)_{-s}\,q ,
\end{equation}
the $s=0$ term being $S_d(a)q$ and the others involving $E^+(a)_{-s}q$ of lower weight $k-s$.

\begin{leancode}
-- VOA/Moonshine/GenerationFock.lean
theorem EEVec_expand (α : L.carrier) {w : ℕ} {p : FockN (rank L)}
    (hp : p ∈ gradingWN (rank L) w) (d : ℕ) :
    EEVec L α (d : ℤ) p
      = ∑ r ∈ Finset.range (w + 1), vSchur L α (d + r) * aExpVec L α r p
\end{leancode}

Call $a$ \emph{safe} for $(\gamma,w)$ if $2\langle a,\gamma-a\rangle+w+1\le0$. Take a safe $a$, $d\ge1$, $q$
of weight $k=w-d$, and $c=\gamma-a$. With $n=-d-1-\langle a,c\rangle$, \eqref{ch04:eq:eamode} and
\eqref{ch04:eq:EEexpand} give
\begin{equation*}
  e^a_{(n)}(q\otimes e^c)=\varepsilon(a,c)\sum_{s=0}^kS_{d+s}(a)E^+(a)_{-s}q\otimes e^\gamma,
\end{equation*}
and the safe-step condition $n\ge k+\langle a,c\rangle$ is exactly safety: substituting
$n=-d-1-\langle a,c\rangle$ it reads $-d-1-\langle a,c\rangle\ge k+\langle a,c\rangle$, i.e.\
$2\langle a,c\rangle+d+k+1\le0$, i.e.\ (since $d+k=w$ and $c=\gamma-a$) $2\langle a,\gamma-a\rangle+w+1\le0$; $q\otimes e^c\in G$ by the induction hypothesis
($k<w$). So the left side is in $G$ by Step 1. By a downward induction on $k$ (the terms with $s\ge1$
have the same shape with larger $d$ and smaller $k$) we get
\begin{equation*}
  S_d(a)\,q\otimes e^\gamma\in G\qquad\text{for all safe }a,\ d\ge1,\ q\in(F_{24})_{w-d}.
\end{equation*}

\emph{Step 3: linearisation.} We want a single oscillator $a_0(-d)$ rather than the Schur
polynomial $S_d(a)$. Replace $a$ by $ka_0$ with $a_0\neq0$ and $k\in\N$: since
$\langle ka_0,\gamma-ka_0\rangle=k\langle a_0,\gamma\rangle-k^2\langle a_0,a_0\rangle\to-\infty$, the vector
$ka_0$ is safe for all large $k$. And $S_d(ka_0)$ is a polynomial in $k$,
\begin{equation*}
  S_d(ka_0)=\sum_{i}k^i\,s_i,\qquad s_1=\tfrac1d\,a_0(-d)\chfourone,
\end{equation*}
because $S_d$ is built from $\tfrac1n a_0(-n)$ by the exponential series, and only the term with a
single oscillator is linear in $k$.

\begin{leancode}
-- VOA/Moonshine/GenerationFock.lean
noncomputable def schurPoly (α : L.carrier) : ℕ → Polynomial (FockN (rank L))
  | 0 => 1
  | (N + 1) => ((N + 1 : ℂ))⁻¹ • ∑ m ∈ Finset.range (N + 1),
      Polynomial.C (vVar L α m) * (Polynomial.X * schurPoly α (N - m))
theorem schurPoly_eval (α : L.carrier) (k : ℕ) (N : ℕ) :
    (schurPoly L α N).eval (algebraMap ℂ (FockN (rank L)) (k : ℂ)) = vSchur L (k • α) N
theorem schurPoly_coeff_one (α : L.carrier) (N : ℕ) :
    (schurPoly L α (N + 1)).coeff 1 = ((N + 1 : ℂ))⁻¹ • vVar L α N
theorem mem_of_forall_sum_pow_smul_mem {M : Type*} [AddCommGroup M] [Module ℂ M]
    (G : Submodule ℂ M) (v : ℕ → M) (D K : ℕ)
    (h : ∀ k : ℕ, K ≤ k → ∑ i ∈ Finset.range (D + 1), ((k : ℂ) ^ i) • v i ∈ G) :
    ∀ i ∈ Finset.range (D + 1), v i ∈ G
\end{leancode}

\begin{dictionary}
\lean{schurPoly L α N} is the polynomial in an auxiliary variable $T$ (\lean{Polynomial.X}) with
coefficients in $F_r$ obtained by putting $T$ in front of every $\alpha(-m-1)$ in the Newton
recursion; \lean{schurPoly_eval} says evaluating at $T=k$ gives $S_N(k\alpha)$, and
\lean{schurPoly_coeff_one} identifies the linear coefficient as $\tfrac1{N+1}\alpha(-N-1)\chfourone$.
\lean{mem_of_forall_sum_pow_smul_mem} is the interpolation lemma: if a polynomial with coefficients
$v_i$ in a vector space takes values in a subspace $G$ at infinitely many integers, then every
$v_i$ lies in $G$.
\end{dictionary}

The interpolation lemma (proof: pass to the quotient by $G$ and apply any linear functional; a
scalar polynomial with infinitely many roots is zero) applied to $k\mapsto S_d(ka_0)q\otimes e^\gamma$
gives $a_0(-d)q\otimes e^\gamma\in G$. Letting $a_0$ run over a basis of $\Leech$ --- whose real
coordinates form an invertible matrix, so the $a_0(-d)$ span all oscillators $h_b(-d)$ --- we get
$h_b(-d)q\otimes e^\gamma\in G$ for every $b$, $d\ge1$ and $q$ of weight $w-d$. Every monomial of
weight $w\ge1$ has this form, which completes the induction.\qed

\begin{example}[Weight two by hand]\label{ch04:ex:genwt2}
Here is the argument in its simplest instance, without scaling. Take $a\in\Leech$, $a\neq0$, and
$n=\langle a,a\rangle-3$. Then
\begin{equation*}
  x_{a\,(n)}x_a=e^a_{(n)}e^{-a}+e^{-a}_{(n)}e^{a}+e^a_{(n)}e^a+e^{-a}_{(n)}e^{-a}.
\end{equation*}
By \eqref{ch04:eq:eamode}, $e^a_{(n)}e^{-a}=\varepsilon(a,-a)EE(a)_{2}\chfourone\otimes e^0=\varepsilon(a,a)S_2(a)$,
and similarly $e^{-a}_{(n)}e^a=\varepsilon(a,a)S_2(-a)$. The last two terms involve $EE(\pm a)_{2-2\langle a,a\rangle}$,
which kills $\chfourone$ since $2-2\langle a,a\rangle<0$. Since $S_2(\pm a)=\tfrac12(a(-1)^2\pm a(-2))\chfourone$,
\begin{equation*}
  x_{a\,(\langle a,a\rangle-3)}\,x_a=\varepsilon(a,a)\,a(-1)^2\chfourone\otimes e^0 .
\end{equation*}
So every $a(-1)^2\chfourone$ lies in the subalgebra generated by the $x_\beta$, and by polarization
$(a+b)(-1)^2-a(-1)^2-b(-1)^2=2a(-1)b(-1)$, so do all quadratic states $h_a(-1)h_b(-1)\chfourone$
(the $a(-1)$ for $a$ in a lattice basis span $\operatorname{span}\{h_b(-1)\}$). Together with the
$x_\beta$ of norm $4$, this is all of $(V_\Leech^+)_2$ (Example~\ref{ch04:ex:wt2}). The general proof
is this computation, organised: polarization is replaced by the interpolation in $k$, which also
handles the higher Schur polynomials.
\end{example}

\begin{remark}[Why no spherical-design input is needed]\label{ch04:rem:design}
Sharper statements are available in the literature: $V_\Leech^+$ is generated by its weight-$2$
space (the Griess algebra), that is, by the quadratic states and the $x_\beta$ with $\beta$ of norm $4$
only. Proving that version from the $196560$ minimal vectors alone requires averaging identities over
the shell $\Leech_4$ --- the fact that the minimal vectors form a spherical design, which makes sums
such as $\sum_{\beta\in\Leech_4}\beta(-1)^2$ proportional to the Casimir $\sum_bh_b(-1)^2$ --- to manufacture
the Fock states. The library proves a weaker statement that suffices for every use: generation by the
$x_\beta$ for \emph{all} $\beta$. Because all multiples $ka_0$ are allowed, the scaling argument of Step 3
replaces any averaging over a shell, and no property of $\Leech$ beyond evenness and
positive-definiteness is used. The cost is nil, since both consumers (\lean{plusForm_fhlInv} below
and \lean{TPForm_fhlInv} in Chapter~\ref{ch:09}) verify their compatibility at $x_\beta$ uniformly
in $\beta$. The design-based route was not attempted in Lean; this remark records why it was not needed.
\end{remark}

\begin{status}
\lean{leechPlus_generated}: proved, unconditional; \texttt{\#print axioms}
\texttt{[propext, Classical.choice, Quot.sound]}. Files \leanpath{VOA/Moonshine/Generation.lean}
and \leanpath{GenerationFock.lean}. (When this chapter was first written, the docstring of \leanpath{Generation.lean} said the file
was ``not wired into the \texttt{VOA/Moonshine.lean} umbrella''; that note was already stale --- the umbrella
imported it. The docstring has since been corrected and now says the file is imported by the umbrella.)
\end{status}
%======================================================================
\section{The invariant form: nondegeneracy of the canonical form on \texorpdfstring{$V_\Lambda^+$}{V+}}\label{ch04:sec:form}

\subsection{What has to be proved, and why}

Every VOA $W$ with $W_0=\C\chfourone$ and $L_1W_1=0$ carries a unique \emph{invariant bilinear form}
normalized by $(\chfourone|\chfourone)=1$ (Li's theorem; in the library, Theorem~\ref{l04:thm:canonical}). For
$V_\Leech^+$ this canonical form is used in Chapter~\ref{ch:09} to define the operator $\overline W$ of
twisted$\times$twisted$\to$untwisted type as the \emph{adjoint} of an operator $W$. Such an adjoint
exists and is unique only if the form is nondegenerate on every weight space. So we need:

\begin{center}
\emph{the canonical invariant form of $V_\Leech^+$ is nondegenerate on each $(V_\Leech^+)_k$.}
\end{center}

The textbook proof is: $V_\Leech^+$ is simple, and the invariant form of a simple VOA is nondegenerate
(its radical is an ideal). The library has the second half (\lean{invariantForm_nondegenerate_FHL})
but not simplicity of $V_\Leech^+$, which is a substantial theorem (Dong--Griess type). The route
taken instead is explicit: write down a form on all of $V_\Leech$, check by hand that it is
nondegenerate, $\theta$-invariant and invariant at the generators $x_\beta$, restrict it to
$V_\Leech^+$, use Theorem~\ref{ch04:thm:gen} to get full invariance, and conclude by uniqueness that
it \emph{is} the canonical form.

\subsection{Invariant forms in the library}

\begin{definition}[Invariant bilinear form]\label{ch04:def:invform}
Let $W$ be a VOA. A bilinear form $(\cdot|\cdot)$ on $W$ is \emph{(fully) invariant} if for every
homogeneous $u\in W_{\mathrm{wt}}$, all $n\in\Z$ and $a,b\in W$,
\begin{equation}\label{ch04:eq:FHL}
  (u_{(n)}a\,|\,b)=(-1)^{\mathrm{wt}}\sum_{j=0}^{\mathrm{wt}}\frac{1}{j!}
  \big(a\,\big|\,(L_1^ju)_{(2\mathrm{wt}-n-2-j)}\,b\big),
\end{equation}
the mode form of $(Y(u,z)a|b)=(a|Y(e^{zL_1}(-z^{-2})^{L_0}u,z^{-1})b)$ (Frenkel--Huang--Lepowsky).
For a \emph{primary} $u$ ($L_1u=0$) only $j=0$ survives:
$(u_{(n)}a|b)=(-1)^{\mathrm{wt}}(a|u_{(2\mathrm{wt}-n-2)}b)$. The form is \emph{normalized} if
$(\chfourone|\chfourone)=1$. The same formula, with $u_{(n)}$ the action on a module $M$ and the
form on $M$, defines invariance of a form on a module at $u$.
\end{definition}

\begin{leancode}
-- VOA/Core/BilinearForm.lean
structure InvariantForm (W : VertexOperatorAlgebra 𝒱) where
  form : LinearMap.BilinForm ℂ 𝒱
  virasoro_adj : ∀ (n : ℤ) (a b : 𝒱), form (W.L n a) b = form a (W.L (-n) b)

def Normalized (B : InvariantForm W) : Prop := B.NormalizedAt 1
def NormalizedAt (B : InvariantForm W) (κ : ℂ) : Prop :=
  B.form W.vacuum W.vacuum = κ

abbrev IsFullyInvariant (B : InvariantForm W) : Prop :=
  ∀ (wt : ℤ) (u : 𝒱), u ∈ W.grading wt → ∀ (n : ℤ) (a b : 𝒱),
    B.form (W.smul u n a) b
      = (-1 : ℂ) ^ wt • ∑ j ∈ Finset.range (wt.toNat + 1), ((j.factorial : ℂ)⁻¹) •
          B.form a (W.smul (((W.L 1) ^ j) u) (2 * wt - n - 2 - (j : ℤ)) b)

def leftRadical (B : InvariantForm W) : Submodule ℂ 𝒱 := LinearMap.ker B.form

-- VOA/Core/InvariantFormFHL2.lean
noncomputable def canonicalInvariantForm (hne : W.vacuum ≠ 0)
    (hqp : ∀ v, v ∈ W.grading 1 → W.L 1 v = 0) : InvariantForm W :=
  canFormFHL hne (canDiagonal_unconditional hne hqp)

theorem invariantForm_unique_FHL (hne : W.vacuum ≠ 0)
    (hqp : ∀ v, v ∈ W.grading 1 → W.L 1 v = 0) {B : InvariantForm W}
    (hB : B.IsFullyInvariant) (hn : B.Normalized) :
    B.form = (canonicalInvariantForm hne hqp).form

-- VOA/Moonshine/InvFormClosure.lean
def FHLInv (N : OrdVModule W M) (B : LinearMap.BilinForm ℂ M) (wt : ℤ) (u : 𝒱) : Prop :=
  ∀ (n : ℤ) (a b : M), B (mode (N.YM u) n a) b
    = (-1 : ℂ) ^ wt • ∑ j ∈ Finset.range (wt.toNat + 1), ((j.factorial : ℂ)⁻¹) •
        B a (mode (N.YM (((W.L 1) ^ j) u)) (2 * wt - n - 2 - (j : ℤ)) b)

-- VOA/Moonshine/WbarRiesz.lean
noncomputable def gradForm (W : VertexOperatorAlgebra 𝒱) (B : LinearMap.BilinForm ℂ 𝒱) (k : ℤ) :
    LinearMap.BilinForm ℂ (W.grading k) :=
  B.compl₁₂ (W.grading k).subtype (W.grading k).subtype
def GradedNondeg (W : VertexOperatorAlgebra 𝒱) (B : LinearMap.BilinForm ℂ 𝒱) : Prop :=
  ∀ k : ℤ, (gradForm W B k).Nondegenerate
\end{leancode}

\begin{dictionary}
\lean{LinearMap.BilinForm ℂ 𝒱} is a bilinear form $\mathcal V\times\mathcal V\to\C$, applied as
\lean{B a b}. An \lean{InvariantForm} in the library is a bilinear form with the Virasoro
adjointness $(L_na|b)=(a|L_{-n}b)$ only; the full condition \eqref{ch04:eq:FHL} is the separate
predicate \lean{IsFullyInvariant}, with $u$ in the first slot, the sum truncated at $j\le\mathrm{wt}$
(\lean{wt.toNat + 1} terms; harmless since $L_1^ju=0$ for $j>\mathrm{wt}$ in a CFT-type VOA), and
\lean{(-1 : ℂ) ^ wt} the sign. \lean{FHLInv N B wt u} is \eqref{ch04:eq:FHL} for one $u$ acting on a
module \lean{N} (\lean{N.YM u} is the module vertex operator); with \lean{N = adjointModule W}, the
VOA as a module over itself (\lean{YM = Y}), it is \eqref{ch04:eq:FHL} on $W$ for that $u$.
\lean{canonicalInvariantForm hne hqp} is the canonical normalized fully invariant form; its
hypotheses are $\chfourone\neq0$ and $L_1W_1=0$. \lean{invariantForm_unique_FHL} is uniqueness: any
fully invariant normalized form equals it (``unconditional'' in the library: no further
hypotheses). \lean{leftRadical} is $\{a:(a|\cdot)=0\}$. \lean{GradedNondeg W B} says the restriction
of $B$ to each weight space $W_k$ (\lean{gradForm}: compose both slots with the inclusion) is
nondegenerate in Mathlib's sense (\lean{Nondegenerate}: only $0$ pairs to zero with everything).
\end{dictionary}

\begin{leancode}
-- VOA/Moonshine/WbarLeech.lean
noncomputable def leechBVForm : InvariantForm leechVOAPlus :=
  canonicalInvariantForm leechVOAPlus_vacuum_ne_zero leechVOAPlus_L1_one
noncomputable abbrev leechBV : LinearMap.BilinForm ℂ SkewVP := leechBVForm.form
theorem leechBV_gradedNondeg_of_isSimple (hsimple : leechVOAPlus.IsSimple) :
    GradedNondeg leechVOAPlus leechBV
\end{leancode}

\begin{dictionary}
\lean{leechBV} is the canonical form of $V_\Leech^+$ (\lean{SkewVP} is another name for the same
type \lean{VPlus}). Its two hypotheses are discharged: $\chfourone\neq0$, and $L_1$ vanishes on
$(V_\Leech^+)_1=0$ (Theorem~\ref{ch04:thm:wtone}). \lean{leechBV_gradedNondeg_of_isSimple} is the
textbook route, \emph{conditional} on simplicity of $V_\Leech^+$, which the library does not prove;
it is superseded by the unconditional Theorem~\ref{ch04:thm:hnd} below and is not used downstream.
\end{dictionary}

\begin{physics}
An invariant bilinear form is the BPZ two-point pairing of states: $(a|b)=\langle a|b\rangle$ computed by conjugating $a$ to the point at infinity, $(Y(u,z)a|b)=(a|Y(e^{zL_1}(-z^{-2})^{L_0}u,z^{-1})b)$. It encodes the two-point function $\langle\phi(\infty)\phi(0)\rangle$ of a chiral field, with the factor $(-1)^{\mathrm{wt}}$ coming from the inversion $z\mapsto-1/z$. Existence and uniqueness (with $(\mathbf 1|\mathbf 1)=1$) mean that two-point functions are fixed by the conformal structure once the vacuum is normalised. Note this form is bilinear, not Hermitian, and is not positive-definite; unitarity (positive-definiteness) holds only on a real form, which is not used here.
\end{physics}

\subsection{The Fock form}

\begin{definition}[The Fock form]\label{ch04:def:uform}
On $F_r$ define the bilinear form $F$ by declaring the monomials orthogonal and
\begin{equation*}
  F(x^d,x^d)=\prod_{(m,a)}\big(-(m+1)\big)^{d_{(m,a)}}\,d_{(m,a)}!\,.
\end{equation*}
\end{definition}

\begin{leancode}
-- VOA/Moonshine/NondegFock.lean
noncomputable def uVar (v : ℕ × Fin r) : ℂ := -((v.1 : ℂ) + 1)
noncomputable def uWt (d : (ℕ × Fin r) →₀ ℕ) : ℂ :=
  d.prod fun v m => uVar v ^ m * (m.factorial : ℂ)
noncomputable def uForm (r : ℕ) : LinearMap.BilinForm ℂ (FockN r) :=
  Finsupp.lsum ℂ fun d => LinearMap.smulRight LinearMap.id (uWt d • lcoeff ℂ d)

theorem uForm_symm (p q : FockN r) : uForm r p q = uForm r q p
theorem uForm_oscN (b : Fin r) (n : ℤ) (p q : FockN r) :
    uForm r (oscN r 1 b n p) q = -uForm r p (oscN r 1 b (-n) q)
theorem uForm_sigmaN (p q : FockN r) : uForm r (sigmaN r p) (sigmaN r q) = uForm r p q
theorem uWt_ne_zero (d : (ℕ × Fin r) →₀ ℕ) : uWt d ≠ 0
theorem uForm_EEVec (γ : L.carrier) (t : ℤ) (y y' : FockN (rank L)) :
    uForm (rank L) (EEVec L γ t y) y' = uForm (rank L) y (EEVec L γ (-t) y')
\end{leancode}

\begin{dictionary}
\lean{uVar (m,a)} $=-(m+1)$ and \lean{uWt d} $=\prod(-(m+1))^{d}\,d!$ (\lean{d.prod} is the product
over the support of $d$). \lean{uForm} is defined on a polynomial $p=\sum_dp_dx^d$ in the first slot
by $F(p,q)=\sum_dp_d\,\mathrm{uWt}(d)\,q_d$ (\lean{lcoeff ℂ d} extracts the coefficient $q_d$): this is
Definition~\ref{ch04:def:uform}. The four theorems say: $F$ is symmetric; $h_b(n)^t=-h_b(-n)$ for
every $n\in\Z$ (the adjoint for $F$); $\sigma$ is an isometry; and every monomial has nonzero norm.
The last line is the transpose of the dressing: $F(EE(\gamma)_ty,y')=F(y,EE(\gamma)_{-t}y')$, i.e.\
$EE(\gamma,z)^t=EE(\gamma,z^{-1})$.
\end{dictionary}

\begin{physics}
The Fock form is the two-point function of the oscillator states: it makes $\alpha^a_{-n}$ and $\alpha^a_n$ conjugate up to the BPZ sign, $(\alpha_n)^t=-\alpha_{-n}$, while a positive Hermitian inner product uses $\alpha_n^\dagger=\alpha_{-n}$. The two differ by the sign $(-1)^{\mathrm{wt}}$ for a weight-$1$ current (the BPZ sign), and each oscillator contributes $-(m+1)$ to the norm. The form is nondegenerate on the Fock space because there are no null states among the oscillator excitations of free bosons.
\end{physics}

\paragraph{Why the minus sign.} We derive the norms from \eqref{ch04:eq:FHL} in three steps.
\emph{(1) Specialize the formula.} Take $u=h_b(-1)\chfourone$. It has weight $\mathrm{wt}=1$, and it
is primary: $L_1u$ would have weight $0$ and charge $0$ and be odd under $\sigma$, but the only
weight-$0$ state is $\chfourone$, which is even; so $L_1u=0$ and only $j=0$ survives in
\eqref{ch04:eq:FHL}. Its field is $Y(u,z)=h_b(z)=\sum_nh_b(n)z^{-n-1}$, so $u_{(n)}=h_b(n)$. With
$\mathrm{wt}=1$ the index on the right is $2\cdot1-n-2=-n$, and \eqref{ch04:eq:FHL} becomes
\begin{equation*}
  (h_b(n)a\,|\,b)=(-1)^1\,(a\,|\,h_b(-n)b)=-(a\,|\,h_b(-n)b)\qquad\text{for all }n,a,b .
\end{equation*}
\emph{(2) Evaluate on the simplest vectors.} Take $n=-1$, $a=\chfourone$, $b=h_b(-1)\chfourone$. Then
$h_b(-1)a=h_b(-1)\chfourone$ and $h_b(1)b=h_b(1)h_b(-1)\chfourone=[h_b(1),h_b(-1)]\chfourone=\chfourone$, so
\begin{equation*}
  (h_b(-1)\chfourone\,|\,h_b(-1)\chfourone)=-(\chfourone\,|\,\chfourone)=-1 .
\end{equation*}
\emph{(3) Iterate.} For a general monomial, move one creation operator $h_b(-m-1)$ from the left
slot to the right; it becomes $-h_b(m+1)=-(m+1)\partial/\partial x_{(m,b)}$, and differentiating
$x_{(m,b)}^{k}$ gives $k\,x_{(m,b)}^{k-1}$. Each removed oscillator therefore contributes a factor
$-(m+1)$ and the multiplicities contribute $d_{(m,b)}!$, which is the norm formula of
Definition~\ref{ch04:def:uform}; monomials with different exponents pair to $0$ because the
derivatives run out. So the form is forced by invariance, minus sign included. Physically: the BPZ conjugation of a weight-$h$
quasi-primary carries the sign $(-1)^h$, so for the current $J=i\partial X$ one gets $J_n^t=-J_{-n}$,
whereas the \emph{Hermitian} (unitary) adjoint is $J_n^\dagger=J_{-n}$. The library's invariant form is
the bilinear (BPZ) one.

\begin{warning}
The ``positive'' Fock form, $h(n)^t=+h(-n)$ (the restriction to real vectors of the familiar
positive-definite inner product of the Fock space), is \emph{not} invariant in the sense of
\eqref{ch04:eq:FHL}: it violates the weight-$1$ case just computed. The project first met this on the
twisted sector, where the sign was checked by computer algebra (the positive form fails, the
negative one works); the same convention is forced here. A form chosen for convenience with the
wrong sign would have made Theorem~\ref{ch04:thm:hnd} below false or vacuous.
\end{warning}

\paragraph{Proof of the Fock-form properties.} Symmetry and $\sigma$-invariance are clear from the
diagonal definition (monomials are $\sigma$-eigenvectors with eigenvalue $\pm1$). For $n=-(m+1)<0$:
$F(x_{(m,b)}x^d,x^{d+e_{(m,b)}})=-(m+1)(d_{(m,b)}+1)F(x^d,x^d)$, while
$h_b(m+1)x^{d+e_{(m,b)}}=(m+1)(d_{(m,b)}+1)x^d$; this is $F(h_b(n)p,q)=-F(p,h_b(-n)q)$ on monomials,
and the case $n>0$ follows by symmetry. For the dressing: $EE(\gamma,z)$ is determined by its
commutator with the oscillators, $[h_b(n),EE(\gamma)_t]=\gamma_b\,EE(\gamma)_{t-n}$ ($n\neq0$), and by its
values on the vacuum, $EE(\gamma)_t\chfourone=S_t(\gamma)$ (or $0$ for $t<0$). Transposing the commutator with
$h_b(n)^t=-h_b(-n)$ shows that $EE(\gamma)_{-t}$ satisfies the transposed relation, and the vacuum
values match since $F(S_t(\gamma),1)=\delta_{t,0}$. Induction on monomials gives the identity.\qed

\begin{status}
\lean{uForm_symm}, \lean{uForm_oscN}, \lean{uForm_sigmaN}, \lean{uWt_ne_zero}, \lean{uForm_EEVec}:
proved, no hypotheses (\leanpath{VOA/Moonshine/NondegFock.lean}).
\end{status}

\subsection{The charge sign and the lattice form}

To extend $F$ to $V_L$ we pair the charge sector $\alpha$ with $-\alpha$:
$B(p\otimes e^\alpha,q\otimes e^\beta)=\delta_{\alpha+\beta,0}\,c(\alpha)\,F(p,q)$ for some signs $c(\alpha)$.
Which $c$? Test invariance at the primary $u=e^\gamma$ of weight $N=\tfrac12\langle\gamma,\gamma\rangle$:
$(e^\gamma_{(n)}x|y)=(-1)^N(x|e^\gamma_{(2N-n-2)}y)$. With $x=p\otimes e^\beta$, $y=q\otimes e^{-\gamma-\beta}$
and \eqref{ch04:eq:Yealpha}, the dressings on the two sides are matched by \lean{uForm_EEVec}
(the indices are $t$ and $-t$), and what remains is
\begin{equation*}
  \varepsilon(\gamma,\beta)\,c(\gamma+\beta)=(-1)^N\varepsilon(\gamma,-\gamma-\beta)\,c(\beta)
  =(-1)^N\varepsilon(\gamma,\gamma)\varepsilon(\gamma,\beta)\,c(\beta),
\end{equation*}
using bimultiplicativity and Proposition~\ref{ch04:prop:cocneg}. So $c(\gamma+\beta)=(-1)^N\varepsilon(\gamma,\gamma)c(\beta)$,
and with $c(0)=1$ the unique solution is
\begin{equation*}
  c(\alpha)=(-1)^{\langle\alpha,\alpha\rangle/2}\,\varepsilon(\alpha,\alpha).
\end{equation*}

\begin{proposition}[The charge sign]\label{ch04:prop:chsign}
For even $L$, $c$ is a character of $L$ with values in $\{\pm1\}$: $c(0)=1$, $c(\alpha+\beta)=c(\alpha)c(\beta)$,
$c(-\alpha)=c(\alpha)$.
\end{proposition}

\begin{leancode}
-- VOA/Moonshine/NondegLattice.lean
noncomputable def chSign (α : L.carrier) : ℂ :=
  (-1 : ℂ) ^ (LatGen.halfNorm L α) * latticeCocycle L α α
theorem chSign_zero : chSign L 0 = 1
theorem chSign_neg (α : L.carrier) : chSign L (-α) = chSign L α
theorem chSign_add (hEven : L.Even) (α β : L.carrier) :
    chSign L (α + β) = chSign L α * chSign L β
\end{leancode}

\begin{dictionary}
\lean{chSign L α} is $c(\alpha)$, with \lean{halfNorm} as in Definition~\ref{ch04:def:lattice}.
\lean{hEven} is needed in \lean{chSign_add} (through the commutator relation and $\tfrac12\langle\cdot,\cdot\rangle$
being additive up to $\langle\alpha,\beta\rangle$).
\end{dictionary}

\begin{physics}
The charge sign $c(\alpha)=(-1)^{\langle\alpha,\alpha\rangle/2}\varepsilon(\alpha,\alpha)$ is the phase in the two-point function $\langle e^{-\alpha}\,e^{\alpha}\rangle$ of a vertex operator with its conjugate: it comes from the BPZ inversion $(-1)^{\mathrm{wt}}$ for weight $\tfrac12\langle\alpha,\alpha\rangle$ and the Klein-factor sign. The fact that $c$ is a character says the sign conventions compose consistently under fusion of momenta, as charge is conserved in the OPE. Convention bookkeeping; no new dynamical content.
\end{physics}

\paragraph{Proof.} $c(\alpha+\beta)=(-1)^{N_\alpha+N_\beta+\langle\alpha,\beta\rangle}\varepsilon(\alpha,\alpha)
\varepsilon(\alpha,\beta)\varepsilon(\beta,\alpha)\varepsilon(\beta,\beta)=c(\alpha)c(\beta)$ by the commutator relation
$\varepsilon(\alpha,\beta)\varepsilon(\beta,\alpha)=(-1)^{\langle\alpha,\beta\rangle}$. Evenness under $-1$ is
Proposition~\ref{ch04:prop:cocneg}.\qed

\begin{warning}[title={Trap: the charge sign is not constant, in general}]
In the literature one usually normalizes $\varepsilon(\alpha,\alpha)=(-1)^{\langle\alpha,\alpha\rangle/2}$; then $c\equiv1$
and the lattice part of the form is simply $(e^\alpha|e^{-\alpha})=1$. The library's basis cocycle does
not satisfy this normalization in general: in coordinates,
$c(\alpha)=(-1)^{\sum_i\frac12\langle B_i,B_i\rangle\alpha_i}$ (the off-diagonal terms of
$\tfrac12\langle\alpha,\alpha\rangle$ cancel against $c(\alpha,\alpha)$, and $\alpha_i^2\equiv\alpha_i$ mod $2$),
which is the trivial character only if every basis vector has norm divisible by $4$. Since the basis
is chosen by \lean{Module.Free.chooseBasis}, nobody can say which case occurs. A form with constant
$c\equiv1$ was the natural first guess and is \emph{wrong} for this cocycle; the project's status
notes record this correction (``NOT constant, because the library's basis cocycle does not satisfy
$\varepsilon(\gamma,\gamma)=(-1)^{\langle\gamma,\gamma\rangle/2}$''). Defining $c$ by the formula above makes
every statement true for \emph{every} basis, which is the only safe option.
\end{warning}

\begin{definition}[The lattice form]\label{ch04:def:latform}
For even $L$, define the bilinear form $B$ on $V_L$ by
\begin{equation*}
  B(p\otimes e^\alpha,q\otimes e^\beta)=\delta_{\alpha+\beta,0}\;c(\alpha)\;F(p,q).
\end{equation*}
\end{definition}

\begin{leancode}
-- VOA/Moonshine/NondegLattice.lean
noncomputable def chForm : LinearMap.BilinForm ℂ (L.carrier →₀ ℂ) :=
  Finsupp.lsum ℂ fun α => LinearMap.smulRight LinearMap.id (chSign L α • Finsupp.lapply (-α))
noncomputable def latForm : LinearMap.BilinForm ℂ (LatModuleGen L) :=
  (uForm (rank L)).tmul (chForm L)

theorem latForm_state (p q : FockN (rank L)) (α β : L.carrier) :
    latForm L (latStateGen L p α) (latStateGen L q β)
      = (if α + β = 0 then chSign L α else 0) * uForm (rank L) p q
theorem latForm_symm (x y : LatModuleGen L) : latForm L x y = latForm L y x
\end{leancode}

\begin{dictionary}
\lean{chForm} is the form on $\C\{L\}$ with $(e^\alpha,f)=c(\alpha)f(-\alpha)$, i.e.\
$(e^\alpha,e^\beta)=\delta_{\alpha+\beta,0}c(\alpha)$ (\lean{Finsupp.lapply (-α)} evaluates at $-\alpha$).
\lean{BilinForm.tmul} is the tensor product of two bilinear forms. \lean{latForm_state} is
Definition~\ref{ch04:def:latform} on basis vectors, and the form is symmetric (using
$c(-\alpha)=c(\alpha)$).
\end{dictionary}

\begin{physics}
This is the BPZ two-point function on the full torus theory: it pairs momentum $\alpha$ with $-\alpha$ (charge conservation: the two-point function is nonzero only for neutral total charge) and the Fock part by the oscillator form. It is the two-point function of the vertex operators $e^\alpha$ and $e^{-\alpha}$ times the oscillator two-point function.
\end{physics}

\begin{theorem}[Properties of the lattice form]\label{ch04:thm:latform}
Let $L$ be even. Then:
\begin{enumerate}[label=(\alph*)]
\item (invariance at $e^\gamma$) for all $\gamma\in L$, $n\in\Z$ and $x,y\in V_L$, with
  $N=\tfrac12\langle\gamma,\gamma\rangle$,
  \begin{equation*}B(e^\gamma_{(n)}x,y)=(-1)^NB(x,e^\gamma_{(2N-n-2)}y);\end{equation*}
\item (nondegeneracy) if $x\neq0$ there is $y$ with $B(x,y)\neq0$;
\item ($\theta$-invariance) $B(\theta x,\theta y)=B(x,y)$;
\item (orthogonality) distinct weight spaces of $V_L$ are $B$-orthogonal.
\end{enumerate}
\end{theorem}

\begin{leancode}
-- VOA/Moonshine/NondegLattice.lean
theorem latForm_Yeαvec (hEven : L.Even) (γ : L.carrier) (n : ℤ) (x y : LatModuleGen L) :
    latForm L (mode (Yeαvec L γ) n x) y
      = (-1 : ℂ) ^ (LatGen.halfNorm L γ) *
          latForm L x (mode (Yeαvec L γ) (2 * LatGen.halfNorm L γ - n - 2) y)
theorem latForm_nondeg (x : LatModuleGen L) (hx : x ≠ 0) : ∃ y, latForm L x y ≠ 0
theorem latForm_thetaLat (x y : LatModuleGen L) :
    latForm L (thetaLat L x) (thetaLat L y) = latForm L x y
theorem latForm_orth_grading {m n : ℤ} (hmn : m ≠ n) {x y : LatModuleGen L}
    (hx : x ∈ gradingLatGen L m) (hy : y ∈ gradingLatGen L n) : latForm L x y = 0
\end{leancode}

\begin{dictionary}
In (a), \lean{mode (Yeαvec L γ) n} is $e^\gamma_{(n)}$: \lean{Yeαvec L γ} is the vertex operator of
$e^\gamma$ in the lattice VOA (it is \lean{YfullGen} of the empty word, \lean{YfullGen_word} with
$l=[\,]$; for Leech this is recorded as \lean{leechVOA_Y_latChargeGen}). $e^\gamma$ is primary of weight
$N$, so (a) is \eqref{ch04:eq:FHL} at $u=e^\gamma$. (a) needs only evenness (not positive-definiteness);
(b)--(d) need no hypothesis beyond the section variable \lean{L}. (b) is nondegeneracy on the whole
space; since weight spaces are orthogonal (d) and finite-dimensional for positive-definite $L$, this
is equivalent to nondegeneracy on each weight space.
\end{dictionary}

\begin{physics}
(a) is the BPZ-invariance of the two-point function at the vertex operator $e^\gamma$: moving $e^\gamma$ from one slot to the other costs $(-1)^N$ and reverses the mode number. (b) says the two-point function has no null states: every nonzero state has a nonzero overlap with some partner. (c) says reflection is an isometry (the form is $\mathbb Z_2$-invariant), so that $V^+$ and $V^-$ are orthogonal. (d) is orthogonality of different energy levels, as for any two-point function of states of definite conformal weight.
\end{physics}

\paragraph{Proof.} (a) On basis vectors this is the computation that determined $c$, run backwards.
(b) Write $x$ in the basis $x^d\otimes e^\alpha$; pick a pair $(d,\alpha)$ with nonzero coefficient. Then
$B(x,x^d\otimes e^{-\alpha})=\text{coefficient}\cdot c(-\alpha)\cdot\mathrm{uWt}(d)\neq0$, since the basis
vectors $x^{d'}\otimes e^{\alpha'}$ pair with $x^d\otimes e^{-\alpha}$ only for $(d',\alpha')=(d,\alpha)$.
(c) $\sigma$ is an $F$-isometry and $c(-\alpha)=c(\alpha)$. (d) The pairing is between sectors $\alpha$ and
$-\alpha$, which have the same half-norm, and $F$ pairs only equal Fock weights.\qed

\begin{example}[The lattice form in low weight]
In weight $1$: $B(h_a(-1)e^0,h_b(-1)e^0)=-\delta_{ab}$. In weight $2$, the basis of
Example~\ref{ch04:ex:wt2} is orthogonal with
\begin{gather*}
  B(h_a(-2)e^0,h_a(-2)e^0)=-2,\qquad B(h_a(-1)^2e^0,h_a(-1)^2e^0)=(-1)^2\,2!=2,\\
  B(h_a(-1)h_b(-1)e^0,h_a(-1)h_b(-1)e^0)=1\ (a\neq b),\qquad
  B(e^\alpha\pm e^{-\alpha},e^\alpha\pm e^{-\alpha})=\pm2c(\alpha).
\end{gather*}
All are nonzero: the form is nondegenerate on $(V_\Leech)_2$ and, separately, on $(V_\Leech^+)_2$ and
$(V_\Leech^-)_2$. On $(V_\Leech)_2$ the form is indefinite (both $2$ and $-2$ occur), as is to be
expected of a bilinear BPZ form; definiteness is never used.
\end{example}

\begin{status}
\lean{latForm_Yeαvec}, \lean{latForm_nondeg}, \lean{latForm_thetaLat},
\lean{latForm_orth_grading}: proved, unconditional (hypothesis \lean{L.Even} where shown);
\texttt{\#print axioms} for \lean{latForm_nondeg} and \lean{latForm_Yeαvec}:
\texttt{[propext, Classical.choice, Quot.sound]} (\leanpath{VOA/Moonshine/NondegLattice.lean}).
\end{status}

\subsection{Restriction to \texorpdfstring{$V_\Lambda^+$}{V+} and identification with the canonical form}

\begin{definition}\label{ch04:def:plusform}
$B^+$ is the restriction of $B$ (for $L=\Leech$) to $V_\Leech^+$.
\end{definition}

\begin{leancode}
-- VOA/Moonshine/NondegLeech.lean
noncomputable def plusForm : LinearMap.BilinForm ℂ VPlus :=
  (Nondeg.latForm leech).compl₁₂ (orbifold (Subgroup.zpowers leechTheta)).carrier.subtype
    (orbifold (Subgroup.zpowers leechTheta)).carrier.subtype

theorem plusForm_fhlInv_frameX0P (β : leech.carrier) :
    FHLInv (adjointModule leechVOAPlus) plusForm (LatGen.halfNorm leech β) (frameX0P β)
theorem plusForm_fhlInv {wt : ℤ} {u : VPlus} (hu : u ∈ leechVOAPlus.grading wt) :
    FHLInv (adjointModule leechVOAPlus) plusForm wt u

noncomputable def plusInvForm : InvariantForm leechVOAPlus where
  form := plusForm
  virasoro_adj := ...
theorem plusInvForm_isFullyInvariant : plusInvForm.IsFullyInvariant
theorem plusInvForm_normalized : plusInvForm.Normalized

theorem leechBV_eq_plusForm : leechBV = plusForm :=
  (invariantForm_unique_FHL leechVOAPlus_vacuum_ne_zero leechVOAPlus_L1_one
    plusInvForm_isFullyInvariant plusInvForm_normalized).symm
\end{leancode}

\begin{dictionary}
\lean{compl₁₂ f g} composes the two slots of a form with linear maps, here both the inclusion
$V_\Leech^+\hookrightarrow V_\Leech$: so \lean{plusForm} is $B^+$. \lean{plusForm_fhlInv_frameX0P}
is \eqref{ch04:eq:FHL} for $u=x_\beta$ (of weight $\tfrac12\langle\beta,\beta\rangle$) on $V_\Leech^+$;
\lean{plusForm_fhlInv} is \eqref{ch04:eq:FHL} for every homogeneous $u\in V_\Leech^+$.
\lean{plusInvForm} packages $B^+$ as an \lean{InvariantForm} (the Virasoro field, elided as
``\lean{...}'', is a proof); it is fully invariant and normalized. The one-line proof of
\lean{leechBV_eq_plusForm} is shown because it is the argument: uniqueness of the normalized fully
invariant form. Note that this identifies the \emph{abstractly defined} canonical form \lean{leechBV}
with the explicit $B^+$; nothing about \lean{leechBV} is assumed.
\end{dictionary}

\begin{physics}
$B^+$ is the BPZ two-point pairing of the orbifold's chiral algebra, restricted from the parent torus theory to the $\theta$-even states. It is the natural inner product on the untwisted sector: correlators in the orbifold of $\mathbb Z_2$-even operators are the same as in the parent theory, so the two-point functions are inherited. The proof that this restricted form is the canonical one uses the generation theorem: invariance at the cosines $x_\beta$ propagates to the whole algebra by the OPE.
\end{physics}

\paragraph{Proof.} \emph{Invariance at $x_\beta$.} $x_\beta=e^\beta+e^{-\beta}$ is primary (it is a sum
of primaries of the same weight), so \eqref{ch04:eq:FHL} at $x_\beta$ is the sum of
Theorem~\ref{ch04:thm:latform}(a) for $\gamma=\pm\beta$.

\emph{Invariance for all $u$.} Let $I\subseteq V_\Leech^+$ be the span of the homogeneous $u$ for which
\eqref{ch04:eq:FHL} holds. It contains $\chfourone$, and it is closed under $n$-th products: this is
the classical argument of Frenkel--Huang--Lepowsky, which derives invariance at $u_{(m)}v$ from
invariance at $u$ and $v$ via the Jacobi identity (the adjoint operator $Y(e^{zL_1}(-z^{-2})^{L_0}u,z^{-1})$
satisfies the opposite Jacobi identity), using that distinct weight spaces are orthogonal
(Theorem~\ref{ch04:thm:latform}(d)). It is proved generically in the library
(\lean{fhlInv_of_generates}, Chapter~\ref{ch:09}). By Theorem~\ref{ch04:thm:gen}, $I=V_\Leech^+$.

\emph{Virasoro adjointness} is the slice $u=\omega$ of full invariance: $\omega$ has weight $2$ and
$L_1\omega\in(V_\Leech^+)_1=0$, so \eqref{ch04:eq:FHL} reads $(L_na|b)=(a|L_{-n}b)$.
\emph{Normalization}: $B(e^0,e^0)=c(0)F(1,1)=1$. \emph{Identification}: uniqueness.\qed

\begin{theorem}[Nondegeneracy of the canonical form of $V_\Leech^+$]\label{ch04:thm:hnd}
The canonical invariant form of $V_\Leech^+$ has zero radical, and is nondegenerate on every weight
space $(V_\Leech^+)_k$.
\end{theorem}

\begin{leancode}
-- VOA/Moonshine/NondegLeech.lean
theorem plusForm_nondeg (x : VPlus) (hx : x ≠ 0) : ∃ y : VPlus, plusForm x y ≠ 0
theorem leechBVForm_leftRadical : InvariantForm.leftRadical leechBVForm = ⊥
theorem leechBV_gradedNondeg : GradedNondeg leechVOAPlus leechBV
\end{leancode}

\begin{dictionary}
\lean{leechBVForm} is the canonical form of $V_\Leech^+$ (\lean{leechBV} is its underlying bilinear
form); \lean{leftRadical … = ⊥} says its radical is zero, and \lean{leechBV_gradedNondeg} is
nondegeneracy on each weight space (Definition: \lean{GradedNondeg}). No hypotheses. This is
precisely the statement called \lean{hnd} in the project notes, the residual that made the operator
$\overline W$ of Chapter~\ref{ch:09} conditional before it was discharged. How could it have been
vacuous? If the form were identically zero, nondegeneracy would fail outright; if it were a form
chosen by hand rather than the canonical one, it would not be the form used by $\overline W$ ---
\lean{leechBV_eq_plusForm} closes that gap.
\end{dictionary}

\begin{physics}
Nondegeneracy of the two-point pairing means the chiral algebra has no null states: no nonzero operator has vanishing two-point function with all operators of the same weight. In CFT terms this is the absence of a nontrivial ideal of null fields (the algebra is simple-like at the level of the invariant form), which is what is needed to define adjoints and conjugate operators in the orbifold. The proof is the physical argument in disguise: the parent theory is nondegenerate and $\theta$ is an isometry, so the pairing of an invariant state with itself and with the projected partner cannot vanish. Positive-definiteness is not claimed (the form is a BPZ bilinear form, indefinite on the complex space).
\end{physics}

\paragraph{Proof.} Let $x\in V_\Leech^+$, $x\neq0$. By Theorem~\ref{ch04:thm:latform}(b) there is $y\in V_\Leech$
with $B(x,y)\neq0$. Then $Py=\tfrac12(y+\theta y)\in V_\Leech^+$ and, using $\theta x=x$ and
$\theta$-invariance,
\begin{equation*}
  B(x,Py)=\tfrac12\big(B(x,y)+B(\theta x,\theta y)\big)=B(x,y)\neq0 .
\end{equation*}
So $B^+$, which is the canonical form, has zero radical. Since distinct weight spaces are orthogonal
and each is finite-dimensional, a vector of $(V_\Leech^+)_k$ pairing to zero with $(V_\Leech^+)_k$ pairs
to zero with everything; hence nondegeneracy on each weight space.\qed

\begin{remark}[Consequences, deferred]
With Theorem~\ref{ch04:thm:hnd} the adjoint $\overline W$ of the skew intertwiner $W$ exists, and the
library's \lean{leechWbar_adj'}, \lean{leechWbar_comm'}, \lean{leechWbar_local'},
\lean{leechWbar_homCov'} (also in \leanpath{NondegLeech.lean}) hold without hypotheses. Their
statements and meaning belong to Chapter~\ref{ch:09}. Nondegeneracy is also what lets the project
prove $\overline W\neq0$, which excludes the square-zero extension
$V_\Leech^+\ltimes(V_\Leech^T)^+$ --- a VOA with $c=24$, $V_1=0$, $\dim V_2=196884$ and graded
dimension $J$, but not simple and not $\Vnat$.
\end{remark}

\begin{status}
\lean{plusForm_fhlInv}, \lean{leechBV_eq_plusForm}, \lean{plusForm_nondeg},
\lean{leechBVForm_leftRadical}, \lean{leechBV_gradedNondeg}: proved, unconditional. \texttt{\#print
axioms} for \lean{leechBV_eq_plusForm} and \lean{leechBV_gradedNondeg}: \texttt{[propext,
Classical.choice, Quot.sound]}. Simplicity of $V_\Leech^+$ is \emph{not} proved and is not needed.
(When this chapter was first written, the file docstring of \leanpath{NondegLeech.lean} said the file was not yet
imported by the umbrella; that note was already stale --- it was imported. The docstring has since been
corrected.)
\end{status}
%======================================================================
\section{Traps and dead ends of this chapter, collected}\label{ch04:sec:traps}

For the referee's convenience we collect the points where the formalization could have gone wrong,
and what prevents it.
\begin{enumerate}[label=(\arabic*)]
\item \emph{Cocycle convention} (\S\ref{ch04:sec:theta}). The naive lift
  $p\otimes e^\alpha\mapsto\sigma p\otimes e^{-\alpha}$ is an automorphism only because the library's
  cocycle is bimultiplicative; for a general cocycle a correcting function is needed. Certified by the
  proved \lean{thetaLat_mode}.
\item \emph{Triviality of $\theta$} (\S\ref{ch04:sec:aut}). If $\theta$ were the identity, $V_\Leech^+$
  would be $V_\Leech$. Excluded by \lean{leechTheta_ne_one} and $98580\neq196884$.
\item \emph{Dimension statements and \lean{finrank}} (\S\ref{ch04:sec:wttwo}). Mathlib's
  \lean{finrank} is $0$ on infinite-dimensional spaces and \lean{Set.ncard} is $0$ on infinite sets,
  so a careless dimension statement could be ``true'' for the wrong reason. All spaces and sets
  counted here are proved finite, and all counts are positive.
\item \emph{Abstract versus concrete weight spaces.} The numbers are proved for
  \lean{leechVOAPlus.grading k}, the weight spaces of the VOA itself, via the explicit linear
  equivalences \lean{leechVOAPlusGradingTwoEquiv} and \lean{leechVOAPlusGradingEquiv} with the
  concrete subspaces \lean{plusTwo}, \lean{plusK}.
\item \emph{$n$-th product generation is not Zhu generation} (\S\ref{ch04:sec:generation}).
  Theorem~\ref{ch04:thm:gen} must not be used where Zhu generation is needed (Chapter~\ref{ch:12}).
\item \emph{Sign of the Fock form} (\S\ref{ch04:sec:form}). The positive Fock form is not invariant;
  the negative (BPZ) one is.
\item \emph{Constant charge sign} (\S\ref{ch04:sec:form}). The choice $(e^\alpha|e^{-\alpha})=1$ is wrong
  for the library's cocycle; the sign $c(\alpha)=(-1)^{\langle\alpha,\alpha\rangle/2}\varepsilon(\alpha,\alpha)$ is
  right for every basis.
\item \emph{Nondegeneracy via simplicity} (\S\ref{ch04:sec:form}). The textbook route needs
  simplicity of $V_\Leech^+$, not in the library; the conditional
  \lean{leechBV_gradedNondeg_of_isSimple} exists but is superseded by the unconditional
  \lean{leechBV_gradedNondeg}.
\item \emph{Stale docstrings (since corrected).} When this chapter was first written, the module docstrings of
  \leanpath{Generation.lean} and \leanpath{NondegLeech.lean} said these files were not wired into the
  umbrella, although they were; at a pre-release build commit both docstrings say the files are imported.
\end{enumerate}

%======================================================================
\section{Referee's checklist}\label{ch04:sec:checklist}

All entries are in \leanpath{VOA/Moonshine/}; ``uncond.'' means: no hypothesis that is an unproved
datum, and \texttt{\#print axioms} $=$ \texttt{[propext, Classical.choice, Quot.sound]} where it was
run (marked $\ast$).

{\footnotesize
\begin{longtable}{>{\raggedright\arraybackslash}p{5.7cm}>{\raggedright\arraybackslash}p{4.4cm}>{\raggedright\arraybackslash}p{2.1cm}>{\raggedright\arraybackslash}p{1.9cm}}
\toprule
Lean declaration & intended meaning & hypotheses & status / check\\
\midrule
\endhead
\lean{latticeCocycle_neg_neg} & $\varepsilon(-\alpha,-\beta)=\varepsilon(\alpha,\beta)$ & none & uncond.; Prop.~\ref{ch04:prop:cocneg}\\
\lean{thetaLat}, \lean{thetaLat_tmul} & $\theta(p\otimes e^\alpha)=\sigma p\otimes e^{-\alpha}$, no correcting sign & none & definition; Def.~\ref{ch04:def:theta}\\
\lean{thetaLat_mode}$^\ast$ & $\theta(u_{(n)}v)=(\theta u)_{(n)}(\theta v)$ & $L$ even, pos.-def. & uncond.; Thm.~\ref{ch04:thm:thetamode}\\
\lean{thetaAut}, \lean{thetaAut_sq}$^\ast$ & $\theta\in\Aut(V_L)$, $\theta^2=1$ & $L$ even, pos.-def. & uncond.; Thm.~\ref{ch04:thm:thetaAut}\\
\lean{thetaAut_ne_one}$^\ast$ & $\theta\neq1$ & $L\neq0$ & uncond.; Thm.~\ref{ch04:thm:thetaAut}\\
\lean{leechTheta}, \lean{leechTheta_ne_one}$^\ast$ & $\theta_\Leech\in\Aut V_\Leech$, $\neq1$ & none & uncond.; Thm.~\ref{ch04:thm:thetaAut}\\
\lean{leechVOAPlus} & $V_\Leech^+=V_\Leech^{\langle\theta\rangle}$ as a VOA & none & definition; Def.~\ref{ch04:def:VPlus}\\
\lean{mem_leechVOAPlus_iff}$^\ast$ & $v\in V_\Leech^+\iff\theta v=v$ & none & uncond.; \S\ref{ch04:sec:plus}\\
\lean{leechVOAPlus_centralCharge}$^\ast$ & $c(V_\Leech^+)=24$ & none & uncond.; \S\ref{ch04:sec:plus}\\
\lean{leechVOAPlus_grading_one}$^\ast$ & $(V_\Leech^+)_1=0$ & none & uncond.; Thm.~\ref{ch04:thm:wtone}\\
\lean{finrank_gradingLatGen_two} & $\dim(V_L)_2=r+\binom{r+1}2+\#L_4$ & even, pos.-def., rootless & uncond.; Thm.~\ref{ch04:thm:wt2rootless}\\
\lean{finrank_plusTwo} & $\dim(V_L^+)_2=\binom{r+1}2+\tfrac12\#L_4$ & even, pos.-def., rootless & uncond.; Thm.~\ref{ch04:thm:plustwo}\\
\lean{leechVOA_finrank_grading_two}$^\ast$ & $\dim(V_\Leech)_2=196884$ & none & uncond.; Thm.~\ref{ch04:thm:wt2rootless}\\
\lean{leechVOAPlus_finrank_grading_two}$^\ast$ & $\dim(V_\Leech^+)_2=98580$ & none & uncond.; Thm.~\ref{ch04:thm:plustwo}\\
\lean{finrank_gradingWN} & $\dim(F_r)_w=P(r,w)$ & none & uncond.; Def.~\ref{ch04:def:fockcount}\\
\lean{hasProd_fockGenFun}$^\ast$ & $\sum P(r,w)q^w=\prod(1-q^n)^{-r}$ & none & uncond.; Thm.~\ref{ch04:thm:genfun}\\
\lean{hasProd_fockSignedGenFun}$^\ast$ & $\sum(P^+-P^-)q^w=\prod(1+q^n)^{-r}$ & none & uncond.; Thm.~\ref{ch04:thm:genfun}\\
\lean{finrank_gradingLatGen_eq_sum_theta} & $\dim(V_L)_k=\sum_n\#L_{2n}P(r,k-n)$ & even, pos.-def. & uncond.; Thm.~\ref{ch04:thm:dimVL}\\
\lean{two_mul_finrank_plusK_eq_signed}$^\ast$ & $2\dim(V_L^+)_k=\dim(V_L)_k+P^{\mathrm s}(r,k)$ & even, pos.-def. & uncond.; Thm.~\ref{ch04:thm:dimVLplus}\\
\lean{leechVOA_finrank_grading_three}$^\ast$ & $\dim(V_\Leech)_3=21493760$ & none & uncond.; Thm.~\ref{ch04:thm:wt3}\\
\lean{leechVOAPlus_finrank_grading_three}$^\ast$ & $\dim(V_\Leech^+)_3=10745856$ & none & uncond.; Thm.~\ref{ch04:thm:wt3}\\
\lean{leechPlus_generated}$^\ast$ & the $x_\beta$ ($\beta\in\Leech$) generate $V_\Leech^+$ under $n$-th products & none & uncond.; Thm.~\ref{ch04:thm:gen}\\
\lean{uForm_oscN}, \lean{uForm_EEVec} & Fock form: $h(n)^t=-h(-n)$, $EE(z)^t=EE(z^{-1})$ & none & uncond.; Def.~\ref{ch04:def:uform}\\
\lean{chSign_add} & $c(\alpha)=(-1)^{N_\alpha}\varepsilon(\alpha,\alpha)$ is a character & $L$ even & uncond.; Prop.~\ref{ch04:prop:chsign}\\
\lean{latForm_Yeαvec}$^\ast$ & invariance of $B$ at every $e^\gamma$ & $L$ even & uncond.; Thm.~\ref{ch04:thm:latform}\\
\lean{latForm_nondeg}$^\ast$ & $B$ nondegenerate on $V_L$ & none & uncond.; Thm.~\ref{ch04:thm:latform}\\
\lean{latForm_thetaLat} & $\theta$ is a $B$-isometry & none & uncond.; Thm.~\ref{ch04:thm:latform}\\
\lean{plusForm_fhlInv} & $B^+$ fully invariant on $V_\Leech^+$ & none & uncond.; \S\ref{ch04:sec:form}\\
\lean{leechBV_eq_plusForm}$^\ast$ & canonical form of $V_\Leech^+$ $=B^+$ & none & uncond.; \S\ref{ch04:sec:form}\\
\lean{leechBV_gradedNondeg}$^\ast$ & canonical form nondegenerate on each $(V_\Leech^+)_k$ (``\lean{hnd}'') & none & uncond.; Thm.~\ref{ch04:thm:hnd}\\
\lean{leechBV_gradedNondeg_of_isSimple} & same, from simplicity & $V_\Leech^+$ simple (open) & conditional, superseded\\
\bottomrule
\end{longtable}
}

%======================================================================
\section*{Notes and references}\addcontentsline{toc}{section}{Notes and references}\label{ch04:sec:notes}

The involution $\theta$, the fixed-point algebra $V_\Leech^+$ and the decomposition
$\Vnat=V_\Leech^+\oplus(V_\Leech^T)^+$ are due to I.~Frenkel, J.~Lepowsky and A.~Meurman, \emph{Vertex
Operator Algebras and the Monster}, Pure and Applied Mathematics 134, Academic Press, Boston, 1988;
the lift of $-1$ with its correcting function is \S10.4 there (the book also treats the central
extensions and cocycles, the graded dimensions and the invariant forms used in this chapter). The earlier
announcement is I.~Frenkel, J.~Lepowsky, A.~Meurman, \emph{A natural representation of the
Fischer--Griess Monster with the modular function $J$ as character}, Proc.\ Natl.\ Acad.\ Sci.\ USA
81 (1984) 3256--3260. Lattice vertex operators and their cocycles go back to R.~Borcherds,
\emph{Vertex algebras, Kac--Moody algebras, and the Monster}, Proc.\ Natl.\ Acad.\ Sci.\ USA 83 (1986)
3068--3071, and to the physics of I.~Frenkel and V.~Kac, \emph{Basic representations of affine Lie
algebras and dual resonance models}, Invent.\ Math.\ 62 (1980) 23--66, and G.~Segal,
\emph{Unitary representations of some infinite-dimensional groups}, Comm.\ Math.\ Phys.\ 80 (1981)
301--342.

Invariant bilinear forms on VOAs and the contragredient formula \eqref{ch04:eq:FHL} are from
I.~Frenkel, Y.-Z.~Huang, J.~Lepowsky, \emph{On axiomatic approaches to vertex operator algebras and
modules}, Mem.\ Amer.\ Math.\ Soc.\ 104 (1993), no.~494; uniqueness and the criterion $L_1V_1=0$ are
from H.~Li, \emph{Symmetric invariant bilinear forms on vertex operator algebras}, J.\ Pure Appl.\
Algebra 96 (1994) 279--297. The textbook treatment followed by the library is T.~Gannon,
\emph{Moonshine beyond the Monster}, Cambridge University Press, 2006, Chapter~5 (lattice VOAs in
\S5.2--5.3, the involution in \S5.3.6). For the structure of $V_L^+$ (simplicity, generation by low
weights, representations) see C.~Dong and R.~Griess, \emph{Rank one lattice type vertex operator
algebras and their automorphism groups}, J.\ Algebra 208 (1998) 262--275; C.~Dong and K.~Nagatomo,
\emph{Representations of vertex operator algebra $V_L^+$ for rank one lattice $L$}, Comm.\ Math.\
Phys.\ 202 (1999) 169--195; and T.~Abe, C.~Dong, H.~Li, \emph{Fusion rules for the vertex operator
algebras $M(1)^+$ and $V_L^+$}, Comm.\ Math.\ Phys.\ 253 (2005) 171--219 (arXiv:math/0310425). The
orbifold viewpoint in physics is L.~Dixon, J.~Harvey, C.~Vafa, E.~Witten, \emph{Strings on orbifolds},
Nucl.\ Phys.\ B 261 (1985) 678--686, and for $\Vnat$ L.~Dixon, P.~Ginsparg, J.~Harvey, \emph{Beauty
and the beast: superconformal symmetry in a Monster module}, Comm.\ Math.\ Phys.\ 119 (1988)
221--241. The spherical-design property of the Leech minimal vectors mentioned in
Remark~\ref{ch04:rem:design} is in J.~H.~Conway and N.~J.~A.~Sloane, \emph{Sphere Packings, Lattices
and Groups}, 3rd ed., Springer, 1999, and B.~Venkov's chapter there.

The Lean material of this chapter was written in the LeanVOA project (repository maintainers
F.~K.~Popov and L.~Rastelli), files \leanpath{VOA/Moonshine/LatticeInvolution.lean} through
\leanpath{NondegLeech.lean}; the route choices recorded here (no correcting sign, the counting trick
instead of traces, generation by all $x_\beta$ via scaling, and the explicit lattice form in place of
simplicity) are documented in the project's
\texttt{Notes/MonsterRoadmap.md} and the ``Status'' sections of \texttt{Notes/M8GluingPlan.md}.
